% This file is meant to be included from `docs/spec.tex` (as an appendix) via `% This file is meant to be included from `docs/spec.tex` (as an appendix) via `% This file is meant to be included from `docs/spec.tex` (as an appendix) via `% This file is meant to be included from `docs/spec.tex` (as an appendix) via `\input{study}`.
% It intentionally contains no `\documentclass`, no preamble, and no `\begin{document}`.

\section{Scope}
We treat \textbf{planning} as the strategic placement and configuration of available
infrastructure under constraints such as interference, and \textbf{dimensioning} as the
minimum resource sizing required to satisfy per-user QoS targets under stochastic
load and resource outages.
Our three core problems of interest are interference, per-user QoS/outage, and route existence.

\section{Common Model Assumptions}
\begin{itemize}
    \item \textbf{Users} form a spatial point process on the Earth surface, within
    a latitude band.
    \item \textbf{Satellites} form a spatial point process on an orbital shell,
    gated by visibility.
    \item \textbf{Radio model} uses a received-power formulation
    $P(x,y)=P_t G_{xy}\|x-y\|^{-\gamma}$ with Rayleigh fading
    $G_{xy}\sim \mathrm{Exp}(1)$ and SINR-based capacity $C=\mathcal{W}\log_2(1+\mathrm{SINR})$.
    This follows standard stochastic-geometry modeling conventions \cite{gomez2018stochastic}.
\end{itemize}

Definitions: $P_t$ is per-satellite transmit power (W); $G_{xy}$ is a unitless
small-scale fading power gain for the link $x\leftrightarrow y$; $\|x-y\|$ is
3D link distance (m); $\gamma$ is the path-loss exponent; $\mathcal{W}$ is bandwidth (Hz);
$\mathrm{SINR}=S/(I+N)$ with $S$ the received signal power, $I$ aggregate
interference, and $N=N_0\mathcal{W}$ the noise power with $N_0$ noise spectral density (W/Hz).

\section{Planning: Interference as the Primary Objective}
We view planning as choosing configuration variables (e.g., satellite density,
altitude, power, bandwidth, and visibility constraints) so that variables of interest are
controlled (primarily intereference) while maintaining acceptable service quality.
\begin{itemize}
    \item \textbf{Interference metric}: for user at $x$ with serving satellite
    $y_0$, define $I(x)=\sum_{y\in \Phi_b\setminus\{y_0\}} P(x,y)$, that is, the aggregate power
    received from all visible satellites except the serving one.
    \item \textbf{Planning target}: enforce a region-level constraint such as
    $\mathbb{P}(\mathrm{SIR}(x)\ge\beta)\ge p_{\mathrm{sir}}$ or
    $\mathbb{E}[I(x)]\le I_{\max}$ for users in the band.
    \item \textbf{Equilibrium framing}: treat operators or satellites as agents
    whose actions affect aggregate interference, then study equilibria under
    congestion penalties.
\end{itemize}

\section{Dimensioning: Per-User QoS Under Network Outage}
Dimensioning seeks the minimum satellite intensity $N$ such that per-user QoS
targets are met with high confidence, even when the network is resource-limited.
\begin{itemize}
    \item \textbf{Per-user outage}: $O_i=1$ if a user is unserved or if
    throughput $C_i < R_{\min}$; otherwise $O_i=0$.
    \item \textbf{Outage rate}: $o = \frac{1}{|\mathcal{U}|}\sum_{i\in\mathcal{U}} O_i$.
    \item \textbf{Dimensioning target}: find the smallest $N$ such that
    $\mathbb{P}(o \le o_{\max}) \ge p_{\mathrm{target}}$ under the stated priors.
\end{itemize}

\section{Routing: Route Existence Under Time-Varying Connectivity}
We study whether an end-to-end route exists between endpoints given a time-varying visibility/connectivity graph
(ground--satellite and inter-satellite links), reporting $\mathbb{P}(\text{route exists})$ (or a certified lower bound).

\section{Planned Experiments}
\begin{itemize}
    \item Vary satellite intensity and compute interference distributions for
    latitude bands to identify feasible planning regimes.
    \item Measure per-user outage and certify feasibility over Monte Carlo trials,
    reporting the minimum $N$ that meets $(o_{\max}, p_{\mathrm{target}})$.
    \item Measure route existence (as a graph-connectivity event) under representative visibility/connectivity snapshots.
\end{itemize}



\section{Theoretical Structure}

\subsection{Topological and Metric Spaces}

\paragraph{Motivation.}
Many sets we work with in applications (such as subsets of $\mathbb{R}^d$) come with a natural notion of
``closeness'' and ``continuity''. This structure is what lets us talk about open sets, continuous costs, and Borel sets,
which in turn are the standard measurable sets used for probability measures.

\paragraph{Intuition.}
A \emph{topology} tells us which subsets of a set should be considered ``open'', generalising open intervals in
$\mathbb{R}$ and open balls in $\mathbb{R}^d$. A \emph{metric} is a concrete way to specify closeness by giving distances.
Compactness is a finiteness condition (no ``escape to infinity'') that makes many existence arguments behave well.

\paragraph{Definition.}
A \textbf{topological space} is a pair $(X,\tau)$ where $\tau\subseteq 2^X$ is a collection of subsets of $X$ (called
\emph{open sets}) such that:
$$
\emptyset\in\tau,\quad X\in\tau,\quad
\bigcup_{i\in I} U_i \in\tau\ \text{for any family }(U_i)_{i\in I}\subseteq\tau,\quad
U\cap V\in\tau\ \text{for all }U,V\in\tau.
$$

A \textbf{metric space} is a pair $(X,d)$ where $d:X\times X\to[0,\infty)$ satisfies (i) $d(x,y)=0\Leftrightarrow x=y$,
(ii) symmetry, and (iii) the triangle inequality. Every metric $d$ induces a topology by taking open sets to be unions of
open balls $B_r(x)=\{z\in X:d(z,x)<r\}$.

A metric space $(X,d)$ is \textbf{compact} if every open cover of $X$ has a finite subcover (equivalently in metric
spaces, every sequence has a convergent subsequence). A \textbf{compact metric space} is a metric space that is compact.

\paragraph{Example.}
$X=[0,1]\subset\mathbb{R}$ with the usual distance $d(x,y)=|x-y|$ is a compact metric space.
More generally, every closed and bounded subset of $\mathbb{R}^d$ with the Euclidean distance is compact.

\subsection{Measurability}

\paragraph{Motivation.}
We wish to make quantitative statements (such as probabilities of events). In order to do so, we
must ensure that the sets we are making statements about are measurable, that is, we must ensure that 
a simple real number can be assigned to their subsets in a consistent manner that correctly reflects their 
relative sizes (masses). For example, we wish to talk about probabilities of events such as 
``signal to noise and interference is below threshold with probability''.

\paragraph{Intuition.}
How do we create a framework that allows us to measure subsets of a set (e.g. events in the case of probability)? 

\paragraph{Definition.}
A \textbf{measurable space} is a pair $(X,\mathcal{F})$ where $X$ is a set and $\mathcal{F}$ is a collection of subsets
of $X$ (called ``events'') on which we will be allowed to assign sizes.
A \textbf{measure} on $(X,\mathcal{F})$ is a function $\mu:\mathcal{F}\to[0,\infty]$ such that $\mu(\emptyset)=0$ and
for any pairwise disjoint sets $A_1,A_2,\dots \in \mathcal{F}$,
$$
\mu\!\left(\bigcup_{n=1}^\infty A_n\right)=\sum_{n=1}^\infty \mu(A_n).
$$
A \textbf{probability measure} is a measure with $\mu(X)=1$.

\paragraph{Example.}
If $X=\{1,2,\dots,n\}$ is finite, a natural choice is $\mathcal{F}=2^X$ (all subsets), and the uniform probability
measure $\mu(A)=|A|/n$. Then $\mu(A)$ is the probability of the event ``the outcome lies in $A$''.

\subsection{$\sigma$-algebras}

\paragraph{Motivation.}
In probability, we want to manipulate events using the usual set operations: we negate events, take unions of events,
and take limits of event sequences. The collection of ``allowed'' events should therefore be closed under these
operations, so that we do not leave the space of events when doing algebra.

\paragraph{Intuition.}
We do \emph{not} assume every subset of $X$ is measurable: we pick a structured family of subsets that is large enough
for the geometric/physical events we care about, but stable under complements and countable unions.

\paragraph{Definition.}
This is done by equipping these sets (e.g. $X$ the set of all agent types, $Y$ the set of all actions) with a 
$\sigma$-algebra $\mathcal{F}$, which is a collection of subsets defined in such a way that assigning a real number 
to each subset in $\mathcal{F}$ is consistent with the usual rules of set operations (union, intersection, complement) 
and countable limits. In other words, we may add, subtract, and take limits while preserving the consistency of our measure.

A $\sigma$-algebra $\mathcal{F}$ on a set $X$ is a collection of subsets of $X$ such that
$$
X \in \mathcal{F}, \qquad
A \in \mathcal{F} \Rightarrow X \setminus A \in \mathcal{F}, \qquad
A_1,A_2,\dots \in \mathcal{F} \Rightarrow \bigcup_{n=1}^\infty A_n \in \mathcal{F}.
$$
These closure properties ensure that measures are consistent under negation and countable limits.

\paragraph{Example.}
When $X\subset \mathbb{R}^d$, we want
to treat ``geometric'' subsets of $X$ as legitimate events, e.g. open balls, half-spaces,
and threshold sets such as $\{x\in X : f(x)\le a\}$.

We therefore take $\mathcal{F}=\mathcal{B}(X)$, the \textbf{Borel $\sigma$-algebra} on $X$,
defined as the \emph{smallest} $\sigma$-algebra that contains all open subsets of $X$
(equivalently, the $\sigma$-algebra generated by open balls).

Open sets matter here because geometric events (``within radius $r$ of a point'', 
``inside a coverage footprint'') are naturally described using open balls and their unions, 
and we want those events to be measurable by default.

\subsection{Measurable Functions}

\paragraph{Motivation.}
We have considered measuring sets (subsets viewed as events). We now consider measuring functions
that represent physical quantities on $X$ (SINR, received power, interference, etc.), because many events
are defined by thresholding such a function.

\paragraph{Intuition.}
Let $x$ be a random location drawn according to a probability measure $\mu$ on $X$
(for example, $\mu$ could be the distribution of user locations in a latitude band).
If $f:X\to\mathbb{R}$ is a physical quantity (such as SINR), then $f(x)$ is a random real number (a random variable).

Terminology note: in this document, when we say ``distribution'' of $x$ we mean a \emph{probability measure}
$\mu\in\mathcal{P}(X)$, so $\mu(X)=1$. (More general, non-normalised measures also appear in modeling---for example,
Poisson point process \emph{intensity} measures---but probabilities are always defined using a probability measure.)

What we want to quantify are statements like ``subset of locations where the SINR is below the threshold $a$''.
This statement corresponds to the event 
$$
\{x\in X : f(x)\le a\}.
$$
Its probability is the $\mu$-mass of that set:
$$
p = \mu(\{x\in X : f(x)\le a\}).
$$

\paragraph{Definition.}
To ensure this definition is consistent (and to allow integrals like averages), we need the set
$\{x\in X : f(x)\le a\}$ to be an element of $\mathcal{F}$ for every threshold $a$.
This is exactly what it means for $f$ to be measurable.

Formally, we equip the domain with a measurable space $(X,\mathcal{F})$ and the real line with
its Borel $\sigma$-algebra $(\mathbb{R},\mathcal{B}(\mathbb{R}))$.
A function $f:X\to\mathbb{R}$ is measurable if for every Borel set $A\in\mathcal{B}(\mathbb{R})$,
the preimage $f^{-1}(A)\in\mathcal{F}$.
In practice, it is enough to check threshold sets: $f$ is measurable if for every $a\in\mathbb{R}$,
$$
\{x\in X : f(x)\le a\}\in\mathcal{F}.
$$

\paragraph{Example.}
Finally, measurability is what makes expectations meaningful: if $f$ is integrable,
the average value of the physical quantity is $\int_X f(x)\,d\mu(x)$.
When $\mathcal{F}=\mathcal{B}(X)$ and $f$ is continuous (as for most geometric/radio quantities),
$f$ is automatically measurable.

\subsection{Push-Forward Measures}
\paragraph{Motivation.}
The notion of a push-forward measure explains how a measure on $X$ (such as the distribution of locations)
induces a corresponding measure on another space through a function. In our setting this ``other space'' is often
$\mathbb{R}$ (e.g., the real-valued SINR), but the same construction applies to maps $X\to Y$ (as in optimal transport).

\paragraph{Intuition.}
Start with a random location $x\sim \mu$ on $X$. Apply a deterministic map $f$ to obtain a new value $y=f(x)$
(for example a real-valued SINR, or a transported location in $Y$). Even though $f$ is deterministic, the output is
random because the input $x$ is random; the push-forward is simply the distribution of $y$.

\paragraph{Definition.}
Formally, let $(X,\mathcal{F})$ and $(Y,\mathcal{G})$ be measurable spaces, let $\mu$ be a measure on $(X,\mathcal{F})$,
and let $f:X\to Y$ be measurable. The \textbf{push-forward} (also called \textbf{image}) measure $f_{\#}\mu$ on
$(Y,\mathcal{G})$ is defined by
$$
(f_{\#}\mu)(B) := \mu(f^{-1}(B)), \qquad B\in\mathcal{G}.
$$
If $\mu$ is a probability measure, then $f_{\#}\mu$ is again a probability measure.
This says: the mass assigned to a set $B$ of outputs equals the mass of the set of inputs that map into $B$.

\paragraph{Example.}
In the common case $Y=\mathbb{R}$ with $\mathcal{G}=\mathcal{B}(\mathbb{R})$, letting $B=(-\infty,a]$ gives
$$
(f_{\#}\mu)((-\infty,a]) = \mu(\{x\in X : f(x)\le a\}),
$$
which is exactly the probability statement we use for threshold events.

\subsubsection{Push-Forward in the Optimal Transport Problem}
\paragraph{Motivation.}
In optimal transport, we want to \emph{move mass} from a source distribution $\mu$ on a space $X$ (a ``pile'')
to a target distribution $\nu$ on a space $Y$ (a ``hole''), while minimizing a cost (the objective).

\paragraph{Intuition.}
A transport map $T:X\to Y$ tells us where each location $x$ sends its mass. The requirement that ``the pile becomes
the hole'' is exactly that transporting $\mu$ through $T$ yields the target distribution $\nu$.

\paragraph{Definition.}
We say that $T$ \emph{pushes $\mu$ forward to $\nu$} if
$$
T_{\#}\mu=\nu,
$$
meaning that for every measurable set $B\subseteq Y$,
$$
\nu(B)=\mu(T^{-1}(B)).
$$
This is exactly the push-forward definition from the previous subsection, with the specific map $f$ renamed to $T$
(to emphasize that it is a transport map) and the codomain being $Y$ rather than $\mathbb{R}$.

\textbf{Kantorovich formulation.} The Monge formulation above restricts us to a \emph{deterministic} assignment:
each $x$ sends all its mass to a single point $T(x)$. This can be too restrictive (for example, if a point of $\mu$
must supply multiple points in $\nu$) and it leads to a nonconvex optimization problem.

The Kantorovich idea is to relax ``one destination per source point'' into a \emph{transport plan} that can split mass:
instead of choosing a map $T$, we choose a joint measure $\gamma$ on $X\times Y$ whose mass at $(x,y)$ represents
how much mass is transported from $x$ to $y$.

In optimisation language, the \textbf{decision variable} (also called the \textbf{optimisation variable}) is the
mathematical object we are free to choose in order to minimise the objective under constraints.
In the Kantorovich formulation, the decision variable is the joint measure $\gamma$ on $X\times Y$.

To represent a valid transport plan from $\mu$ to $\nu$, $\gamma$ should have marginals $\mu$ and $\nu$.
Equivalently, the constraints that $\gamma$ has marginals equal to $\mu$ and $\nu$ can be written as push-forward relations
$$
(\pi_X)_{\#}\gamma=\mu,\qquad (\pi_Y)_{\#}\gamma=\nu,
$$
where $\pi_X(x,y)=x$ and $\pi_Y(x,y)=y$ are the coordinate projections.
These are exactly push-forwards of the same form $T_{\#}(\cdot)$, with the ``map'' $T$ being instantiated as
$T=\pi_X$ and $T=\pi_Y$, respectively.

To make the origin and endpoint explicit: $\gamma$ is a measure on the \emph{product space} $X\times Y$.
The map $\pi_X:X\times Y\to X$ takes a pair $(x,y)$ and returns its $X$-coordinate $x$, so $(\pi_X)_{\#}\gamma$
is a measure \emph{on $X$}. Likewise, $\pi_Y:X\times Y\to Y$ returns the $Y$-coordinate, so $(\pi_Y)_{\#}\gamma$
is a measure \emph{on $Y$}.

\paragraph{Example.}
Let $X=[0,1]$, $Y=[0,2]$, and let $\mu$ be the uniform distribution on $[0,1]$. Define $T(x)=2x$.
Then $T_{\#}\mu$ is the uniform distribution on $[0,2]$ (the interval is stretched, so density halves), i.e.\ $\nu=T_{\#}\mu$.

Kantorovich-only example (mass splitting). Let $X=\{0\}$ and $Y=\{-1,+1\}$. Let $\mu=\delta_0$ (all mass at $0$) and
let $\nu=\tfrac12\delta_{-1}+\tfrac12\delta_{+1}$. There is no map $T:X\to Y$ such that $T_{\#}\mu=\nu$, because any map
must send all mass to a single point in $Y$. But there is a transport plan $\gamma$ on $X\times Y$ defined by
$$
\gamma(0,-1)=\tfrac12,\qquad \gamma(0,+1)=\tfrac12,
$$
which satisfies $(\pi_X)_{\#}\gamma=\mu$ and $(\pi_Y)_{\#}\gamma=\nu$.

\section{Optimal Transport and Cournot-Nash Equilibria}
This section summarises the background and key ideas in \cite{blanchet2016optimal} in a format aligned with our study.

We are interested in settings with many agents, each attempting to optimise for itself. What sorts of equilibria arise?
What do these equilibria look like? When do they exist, and when are they unique?

For example, consider a large number of users (agents) choosing which satellite (action) to connect to. 
Each user's utility depends on which satellite it chooses, but also on how many other users choose the 
same satellite (congestion externality). We wish to characterise the equilibria of this system, so we may 
understand not only how users behave but also how to design the system (e.g., satellite parameters) to 
achieve desirable outcomes.

Each user is solving an individual optimisation problem, but the problems are coupled through the congestion effects.
We need to characterise jointly the optimisation problem as well as the equilibrium conditions. It would not suffice 
to simply solve each user's optimisation problem in isolation, or to show that an equilibrium exists. It is essential 
to simultaneously characterise individual optimisation strategies and the resulting equilibrium.

In game-theoretic terms, we are interested in \textbf{non-atomic games with externalities}, where each agent is infinitesimally
small (non-atomic) and the cost incurred by each agent depends on the aggregate behaviour of all agents (externalities).

\subsection{Problem Class, Background, and Gap}

\cite{blanchet2016optimal} studies non-atomic games with externalities and gives classes of models where Cournot--Nash
equilibria are not only shown to exist, but can also be characterised in a way that supports uniqueness results and
computation. Their stated gap is that, while abstract existence results for Cournot--Nash equilibria are well known, they
do not typically yield a tractable characterisation (or numerical scheme) for the equilibrium distribution.

\paragraph{Where do externalities and congestion come from?}
In many applications, an agent's cost depends on what \emph{others} do through an aggregate quantity. For instance, in a
wireless association model, if many users choose the same satellite (or the same beam), then service quality degrades;
in road choice, if many drivers choose the same road, travel time increases. These effects are called \emph{externalities}
because an agent's cost depends on the aggregate behaviour of the population, not only on the agent's own action.
The paper models such effects through an action distribution $\nu$ and an externality term $V[\nu](y)$ that modifies the
cost of playing action $y$.

However, once we introduce such externalities (especially congestion-type effects), there is a key technical
obstruction: the dependence of costs on the population distribution can fail to have the regularity properties assumed
by standard equilibrium existence proofs.

\paragraph{What is the ``obstruction''?}
Here, an \emph{obstruction} simply means a technical barrier that prevents us from applying the most common existence
proof techniques ``out of the box''. Concretely, many existence proofs for equilibria rely on a fixed-point argument and
therefore need a basic stability property: if the population action distribution $\nu$ changes a little, then the cost
landscape faced by agents should also change a little. Without such stability, it becomes hard to even define a
well-behaved ``best-response of the population'' map, let alone prove it has a fixed point.

Congestion effects are often \emph{local}: the cost at action $y$ depends on how crowded that action is. A common model is
$V[\nu](y)=f(\nu(y))$, where $\nu(y)$ denotes a density. The difficulty is that densities are fragile: two action
distributions can look almost the same when you only look at coarse averages (e.g.\ averages over regions), yet have very
different pointwise densities because mass can concentrate into narrow spikes or oscillate rapidly. If $V[\nu](y)$ depends
directly on $\nu(y)$, these small ``distribution-level'' changes can produce large changes in the local congestion cost,
so the regularity assumptions behind standard fixed-point existence arguments fail.

To make these statements precise, we first fix the paper's equilibrium notation (Subsection~\ref{sec:cn_equilibrium}) and then explain
the relevant notion of ``two distributions being close'' (weak-$*$ topology) and why local congestion breaks it
(Subsection~\ref{sec:cn_weakstar_gap}).

The key contribution in \cite{blanchet2016optimal} is to exploit an optimal-transport structure (and a variational formulation in potential
cases) to turn equilibrium computation into an optimisation problem over distributions.



\subsection{Equilibrium}\label{sec:cn_equilibrium}
In a finite-player game, one can describe an equilibrium by listing each player's action (or mixed strategy).
In a \emph{non-atomic} game with a continuum of agents, this is not the right object: we do not label agents individually,
and a single agent has negligible influence on the aggregate outcome. Instead, we describe equilibria at the population
level using probability measures.

\subsubsection{Spaces and Probability Measures}
In the paper's notation, $X$ is the \emph{type space} (agent characteristics) and $Y$ is the \emph{action space}.
We write $\mathcal{P}(X)$ for the set of \emph{Borel} probability measures on $X$ and similarly $\mathcal{P}(Y)$.
Here ``Borel'' means the measures are defined on the Borel $\sigma$-algebra generated by the topology on $X$ (and $Y$),
which is the standard choice when $X$ and $Y$ are compact metric spaces. (The ``continuum of agents'' aspect is captured
by describing populations via measures.)
The population of types is given exogenously by a probability measure $\mu\in\mathcal{P}(X)$.

\subsubsection{Equilibrium as a Joint Distribution}
What matters is the population-level assignment of actions to types: for each type $x$, what fraction of agents choose
each action $y$. This assignment is represented by a joint probability measure on $X\times Y$.

An equilibrium is described by a joint probability measure $\gamma\in\mathcal{P}(X\times Y)$, where (informally)
$\gamma(A\times B)$ is the probability that an agent has type in $A\subseteq X$ and chooses an action in $B\subseteq Y$.
The first marginal of $\gamma$ is the induced distribution of types and must equal $\mu$.
The second marginal of $\gamma$ is the induced distribution of actions, denoted $\nu\in\mathcal{P}(Y)$.
Using the coordinate projections $\pi_X(x,y)=x$ and $\pi_Y(x,y)=y$, these marginal constraints can be written as
push-forward relations:
$$
(\pi_X)_{\#}\gamma=\mu,\qquad \nu=(\pi_Y)_{\#}\gamma.
$$
If the equilibrium is \emph{pure}, then there is a map $T:X\to Y$ such that $\gamma$ is carried by the graph of $T$,
i.e.\ $\gamma=(\mathrm{id},T)_{\#}\mu$.

\subsubsection{Costs and Externalities}
Costs are additively separable: an $x$-type agent who plays action $y$ incurs a cost of the form
$$
\Phi(x,y,\nu)=c(x,y)+V[\nu](y),
$$
where $c(x,y)$ is an idiosyncratic term (distance, mismatch, etc.), and the map $\nu\mapsto V[\nu]$ captures
externalities induced by the aggregate action distribution $\nu$.

\subsubsection{Reference Measure, Densities, and Local Congestion}
To speak about ``local density'' at an action $y$, we need $\nu$ to admit a density with respect to some reference
measure on $Y$.
Fix a reference measure $m_0$ on $Y$ and restrict to action distributions $\nu$ that are absolutely continuous with
respect to $m_0$ (written $\nu\ll m_0$). In that case, $\nu$ admits a density $\frac{d\nu}{dm_0}$, and with a slight abuse
of notation we write $\nu(y)$ for this density value at $y$.
Local congestion costs take the form
$$
V[\nu](y)=f(\nu(y)),
$$
so the externality cost at $y$ depends pointwise on how concentrated the action distribution is at $y$.

\subsubsection{Cournot--Nash Equilibrium Condition (Formal Definition)}\label{sec:cn_equilibrium_definition}
\paragraph{Motivation.}
So far, $\gamma$ is just a joint distribution that matches the exogenous type distribution $\mu$ and induces an action
distribution $\nu$. To call $\gamma$ an \emph{equilibrium}, we must encode the behavioural requirement that agents choose
cost-minimising actions \emph{given} the aggregate action distribution $\nu$ that results from everyone's choices.

\paragraph{Intuition.}
Fix a candidate aggregate action distribution $\nu$. Each type $x$ then faces a deterministic optimisation problem:
choose $y\in Y$ to minimise $c(x,y)+V[\nu](y)$. An equilibrium is a self-consistent joint distribution $\gamma$ such that
(i) its second marginal is exactly the $\nu$ used in the optimisation, and (ii) $\gamma$ puts mass only on pairs $(x,y)$
where $y$ is a best response for type $x$.

\paragraph{Definition.}
Let $\gamma\in\mathcal{P}(X\times Y)$ and let $\nu=(\pi_Y)_{\#}\gamma$ be its second marginal.
We say that $\gamma$ is a \textbf{Cournot--Nash equilibrium} (in distributional form) if
\begin{enumerate}[label=(\roman*)]
    \item $(\pi_X)_{\#}\gamma=\mu$ (types are distributed according to $\mu$), and
    \item $\gamma$ is concentrated on best responses given $\nu$, meaning
    $$
    \gamma\bigl(\{(x,y)\in X\times Y:\ c(x,y)+V[\nu](y)=\min_{z\in Y}\bigl(c(x,z)+V[\nu](z)\bigr)\}\bigr)=1.
    $$

    In probabilistic terms: let $(\Omega,\mathcal{F},\mathbb{P})$ be a probability space and let $(X_{\Omega},Y_{\Omega})$
    be an $(X\times Y)$-valued random variable with law $\gamma$ (so $\mathbb{P}[(X_{\Omega},Y_{\Omega})\in A]=\gamma(A)$ for
    every measurable $A\subseteq X\times Y$). Define the event
    $$
    E := \Bigl\{c(X_{\Omega},Y_{\Omega})+V[\nu](Y_{\Omega})=\min_{z\in Y}\bigl(c(X_{\Omega},z)+V[\nu](z)\bigr)\Bigr\}.
    $$
    Then the condition $\gamma(\cdot)=1$ above is exactly the statement $\mathbb{P}(E)=1$: with probability one, the sampled
    action $Y_{\Omega}$ is a best response for the sampled type $X_{\Omega}$ given the aggregate distribution $\nu$.

\end{enumerate}


\paragraph{Intuition.}
This expresses a ``best response for almost every agent'' condition in a non-atomic (continuum) model: we do \emph{not}
quantify over individually labelled agents. Instead, we quantify over the population measure $\gamma$: the set of
non-best-response pairs has $\gamma$-mass $0$, i.e.\ if we pick an agent at random from the population (sample
$(x,y)\sim\gamma$), then with probability $1$ the realised action $y$ is a best response for the realised type $x$.

Equivalently, if we define the best-response set $B_{\nu}\subseteq X\times Y$ by the bracketed condition inside
$\gamma(\cdot)$ above, then $\gamma(B_{\nu})=1$ and the event $E$ is exactly $\{(X_{\Omega},Y_{\Omega})\in B_{\nu}\}$.

\paragraph{Alternative Formulation.}

Equivalently, define the \emph{value function}
$$
\varphi(x):=\min_{z\in Y}\bigl(c(x,z)+V[\nu](z)\bigr),
$$
then the equilibrium condition above can be written as the pair of inequalities
$$
c(x,y)+V[\nu](y)\ge \varphi(x)\quad\text{for all }x\in X\text{ and (where defined) }y\in Y,
$$
with equality holding $\gamma$-almost everywhere (i.e.\ for $\gamma$-almost every pair $(x,y)$).

\paragraph{Example.}
In a satellite-association toy model, let $x$ be a user location and let $y$ index a satellite/beam choice. The term
$c(x,y)$ can represent a geometry-based penalty (e.g.\ path loss), and $V[\nu](y)$ can represent a congestion penalty for
how many users choose satellite/beam $y$ (captured by the aggregate action distribution $\nu$). The equilibrium
condition says that, under the self-consistent distribution of choices, almost every user is assigned only to a choice
that minimises its own cost given the congestion created by all users.


\paragraph{Motivation.}
The equilibrium definition in Subsubsection~\ref{sec:cn_equilibrium_definition} is a self-consistent best-response
condition on a joint distribution $\gamma$.
This definition is conceptually clear, but it does not yet tell us how to \emph{compute} an equilibrium, or how to obtain
uniqueness and stability properties.

\paragraph{Intuition.}
The paper's key step is to identify an important class of models where the equilibrium conditions are exactly the
first-order optimality conditions of a single global objective. In that case, equilibrium computation becomes an
optimisation problem over distributions. Optimal transport enters through the transport cost term, while congestion and
interactions enter through an ``energy'' functional of the action distribution.

\paragraph{Definition.}
In the \emph{potential} (variational) case, the externality map $\nu\mapsto V[\nu]$ is assumed to be the first variation
of some functional $\mathcal{E}$ defined on action distributions:
$$
V[\nu] \;=\; \frac{\delta \mathcal{E}}{\delta \nu}.
$$
Under this assumption, equilibrium action distributions $\nu$ can be characterised as minimisers of an OT-structured
functional of the form
$$
\nu \in \arg\min_{\tilde\nu}\ \mathcal{W}_c(\mu,\tilde\nu) + \mathcal{E}[\tilde\nu],
$$
and an equilibrium coupling $\gamma$ is then obtained as an optimal transport plan between $\mu$ and that minimising
$\nu$.

\paragraph{Example.}
A standard ``good'' potential form used in the paper is
$$
\mathcal{E}[\nu] \;=\; \int_Y F(y,\nu(y))\,dm_0(y) \;+\; \frac12\iint_{Y\times Y} \phi(y,z)\,d\nu(y)\,d\nu(z),
$$
where the first term is a local congestion energy (with $F'(\cdot)=f(\cdot)$) and the second term is a symmetric pairwise
interaction energy (symmetry is what makes it a potential).

\subsection{Why Variational Structure Matters}\label{sec:cn_variational_importance}

\paragraph{Motivation.}
The ability to write equilibrium computation as a global minimisation problem is the reason this approach becomes
practical. Without a variational structure, we typically lose both tractability and the clean uniqueness/stability
statements that come from convexity (or convexity along transport geodesics).

\paragraph{Intuition.}
For general Nash equilibria, there is no reason that ``the equilibrium'' should be the minimiser of a single scalar
objective: different players optimise different objectives, and their incentives can create cycles and rotational
behaviour. A variational (potential) structure is exactly the condition that makes the equilibrium conditions integrable
into a single energy landscape: equilibria are stationary points of that landscape, and in the convex cases they are
global minimisers.

In the Cournot--Nash setting considered by \cite{blanchet2016optimal}, the variational structure is what turns a
definition-by-best-response (Subsubsection~\ref{sec:cn_equilibrium_definition}) into a computable characterisation
(Subsection~\ref{sec:cn_variational_formulation}): solve a minimisation problem for the action distribution $\nu$, then
recover the coupling $\gamma$ using optimal transport.

\paragraph{Definition.}
We will refer to the model as \textbf{variational} (or \textbf{potential}) if there exists a functional
$\mathcal{E}$ on action distributions such that $V[\nu]=\delta\mathcal{E}/\delta\nu$.
In this case, the equilibrium conditions coincide with the first-order optimality conditions of the global objective
$\mathcal{W}_c(\mu,\nu)+\mathcal{E}[\nu]$.
If, in addition, this objective is (geodesically) convex in a suitable sense, then the equilibrium is unique and stable
to perturbations of the inputs $(\mu,c,\mathcal{E})$.

\paragraph{Example.}
Not every game has this structure. In finite games such as Rock--Paper--Scissors or Matching Pennies, incentives create
cyclic dominance, and there is no single potential function whose value strictly improves along unilateral best-response
moves. In continuous-action games, strong asymmetric cross-effects can lead to non-monotone best-response operators (the
system behaves more like a rotation than a gradient flow). The relevance here is conceptual: the Blanchet--Carlier
approach relies on being able to convert equilibrium computation into minimisation, which is precisely what fails when a
game is not variational.

\subsection{Practical Takeaway: Variational Characterisation vs.\ Continuum Complexity}
\label{sec:cn_weakstar_gap}
\label{sec:cn_variational_formulation}
The main technical payoff in \cite{blanchet2016optimal} is that, in the \emph{potential} (variational) case, equilibrium
computation can be reframed as a single global optimisation problem (an optimal-transport assignment term plus an energy
term in the action distribution). This provides a principled \emph{characterisation} of equilibrium and, in favourable
convex cases, supports uniqueness and stable numerical schemes.

For this study, however, pursuing the full continuum route is not an attractive path: the formulation lives on
continuum objects (measures on compact metric spaces, densities, and associated regularity/continuity assumptions), while
our satellite questions are naturally finite and simulator-driven. We therefore retain the equilibrium definition as a
useful theoretical lens, but we do not attempt to implement or rely on the full continuum variational machinery.



% References for citations used here are provided by the parent document (`docs/spec.tex`).
`.
% It intentionally contains no `\documentclass`, no preamble, and no `\begin{document}`.

\section{Scope}
We treat \textbf{planning} as the strategic placement and configuration of available
infrastructure under constraints such as interference, and \textbf{dimensioning} as the
minimum resource sizing required to satisfy per-user QoS targets under stochastic
load and resource outages.
Our three core problems of interest are interference, per-user QoS/outage, and route existence.

\section{Common Model Assumptions}
\begin{itemize}
    \item \textbf{Users} form a spatial point process on the Earth surface, within
    a latitude band.
    \item \textbf{Satellites} form a spatial point process on an orbital shell,
    gated by visibility.
    \item \textbf{Radio model} uses a received-power formulation
    $P(x,y)=P_t G_{xy}\|x-y\|^{-\gamma}$ with Rayleigh fading
    $G_{xy}\sim \mathrm{Exp}(1)$ and SINR-based capacity $C=\mathcal{W}\log_2(1+\mathrm{SINR})$.
    This follows standard stochastic-geometry modeling conventions \cite{gomez2018stochastic}.
\end{itemize}

Definitions: $P_t$ is per-satellite transmit power (W); $G_{xy}$ is a unitless
small-scale fading power gain for the link $x\leftrightarrow y$; $\|x-y\|$ is
3D link distance (m); $\gamma$ is the path-loss exponent; $\mathcal{W}$ is bandwidth (Hz);
$\mathrm{SINR}=S/(I+N)$ with $S$ the received signal power, $I$ aggregate
interference, and $N=N_0\mathcal{W}$ the noise power with $N_0$ noise spectral density (W/Hz).

\section{Planning: Interference as the Primary Objective}
We view planning as choosing configuration variables (e.g., satellite density,
altitude, power, bandwidth, and visibility constraints) so that variables of interest are
controlled (primarily intereference) while maintaining acceptable service quality.
\begin{itemize}
    \item \textbf{Interference metric}: for user at $x$ with serving satellite
    $y_0$, define $I(x)=\sum_{y\in \Phi_b\setminus\{y_0\}} P(x,y)$, that is, the aggregate power
    received from all visible satellites except the serving one.
    \item \textbf{Planning target}: enforce a region-level constraint such as
    $\mathbb{P}(\mathrm{SIR}(x)\ge\beta)\ge p_{\mathrm{sir}}$ or
    $\mathbb{E}[I(x)]\le I_{\max}$ for users in the band.
    \item \textbf{Equilibrium framing}: treat operators or satellites as agents
    whose actions affect aggregate interference, then study equilibria under
    congestion penalties.
\end{itemize}

\section{Dimensioning: Per-User QoS Under Network Outage}
Dimensioning seeks the minimum satellite intensity $N$ such that per-user QoS
targets are met with high confidence, even when the network is resource-limited.
\begin{itemize}
    \item \textbf{Per-user outage}: $O_i=1$ if a user is unserved or if
    throughput $C_i < R_{\min}$; otherwise $O_i=0$.
    \item \textbf{Outage rate}: $o = \frac{1}{|\mathcal{U}|}\sum_{i\in\mathcal{U}} O_i$.
    \item \textbf{Dimensioning target}: find the smallest $N$ such that
    $\mathbb{P}(o \le o_{\max}) \ge p_{\mathrm{target}}$ under the stated priors.
\end{itemize}

\section{Routing: Route Existence Under Time-Varying Connectivity}
We study whether an end-to-end route exists between endpoints given a time-varying visibility/connectivity graph
(ground--satellite and inter-satellite links), reporting $\mathbb{P}(\text{route exists})$ (or a certified lower bound).

\section{Planned Experiments}
\begin{itemize}
    \item Vary satellite intensity and compute interference distributions for
    latitude bands to identify feasible planning regimes.
    \item Measure per-user outage and certify feasibility over Monte Carlo trials,
    reporting the minimum $N$ that meets $(o_{\max}, p_{\mathrm{target}})$.
    \item Measure route existence (as a graph-connectivity event) under representative visibility/connectivity snapshots.
\end{itemize}



\section{Theoretical Structure}

\subsection{Topological and Metric Spaces}

\paragraph{Motivation.}
Many sets we work with in applications (such as subsets of $\mathbb{R}^d$) come with a natural notion of
``closeness'' and ``continuity''. This structure is what lets us talk about open sets, continuous costs, and Borel sets,
which in turn are the standard measurable sets used for probability measures.

\paragraph{Intuition.}
A \emph{topology} tells us which subsets of a set should be considered ``open'', generalising open intervals in
$\mathbb{R}$ and open balls in $\mathbb{R}^d$. A \emph{metric} is a concrete way to specify closeness by giving distances.
Compactness is a finiteness condition (no ``escape to infinity'') that makes many existence arguments behave well.

\paragraph{Definition.}
A \textbf{topological space} is a pair $(X,\tau)$ where $\tau\subseteq 2^X$ is a collection of subsets of $X$ (called
\emph{open sets}) such that:
$$
\emptyset\in\tau,\quad X\in\tau,\quad
\bigcup_{i\in I} U_i \in\tau\ \text{for any family }(U_i)_{i\in I}\subseteq\tau,\quad
U\cap V\in\tau\ \text{for all }U,V\in\tau.
$$

A \textbf{metric space} is a pair $(X,d)$ where $d:X\times X\to[0,\infty)$ satisfies (i) $d(x,y)=0\Leftrightarrow x=y$,
(ii) symmetry, and (iii) the triangle inequality. Every metric $d$ induces a topology by taking open sets to be unions of
open balls $B_r(x)=\{z\in X:d(z,x)<r\}$.

A metric space $(X,d)$ is \textbf{compact} if every open cover of $X$ has a finite subcover (equivalently in metric
spaces, every sequence has a convergent subsequence). A \textbf{compact metric space} is a metric space that is compact.

\paragraph{Example.}
$X=[0,1]\subset\mathbb{R}$ with the usual distance $d(x,y)=|x-y|$ is a compact metric space.
More generally, every closed and bounded subset of $\mathbb{R}^d$ with the Euclidean distance is compact.

\subsection{Measurability}

\paragraph{Motivation.}
We wish to make quantitative statements (such as probabilities of events). In order to do so, we
must ensure that the sets we are making statements about are measurable, that is, we must ensure that 
a simple real number can be assigned to their subsets in a consistent manner that correctly reflects their 
relative sizes (masses). For example, we wish to talk about probabilities of events such as 
``signal to noise and interference is below threshold with probability''.

\paragraph{Intuition.}
How do we create a framework that allows us to measure subsets of a set (e.g. events in the case of probability)? 

\paragraph{Definition.}
A \textbf{measurable space} is a pair $(X,\mathcal{F})$ where $X$ is a set and $\mathcal{F}$ is a collection of subsets
of $X$ (called ``events'') on which we will be allowed to assign sizes.
A \textbf{measure} on $(X,\mathcal{F})$ is a function $\mu:\mathcal{F}\to[0,\infty]$ such that $\mu(\emptyset)=0$ and
for any pairwise disjoint sets $A_1,A_2,\dots \in \mathcal{F}$,
$$
\mu\!\left(\bigcup_{n=1}^\infty A_n\right)=\sum_{n=1}^\infty \mu(A_n).
$$
A \textbf{probability measure} is a measure with $\mu(X)=1$.

\paragraph{Example.}
If $X=\{1,2,\dots,n\}$ is finite, a natural choice is $\mathcal{F}=2^X$ (all subsets), and the uniform probability
measure $\mu(A)=|A|/n$. Then $\mu(A)$ is the probability of the event ``the outcome lies in $A$''.

\subsection{$\sigma$-algebras}

\paragraph{Motivation.}
In probability, we want to manipulate events using the usual set operations: we negate events, take unions of events,
and take limits of event sequences. The collection of ``allowed'' events should therefore be closed under these
operations, so that we do not leave the space of events when doing algebra.

\paragraph{Intuition.}
We do \emph{not} assume every subset of $X$ is measurable: we pick a structured family of subsets that is large enough
for the geometric/physical events we care about, but stable under complements and countable unions.

\paragraph{Definition.}
This is done by equipping these sets (e.g. $X$ the set of all agent types, $Y$ the set of all actions) with a 
$\sigma$-algebra $\mathcal{F}$, which is a collection of subsets defined in such a way that assigning a real number 
to each subset in $\mathcal{F}$ is consistent with the usual rules of set operations (union, intersection, complement) 
and countable limits. In other words, we may add, subtract, and take limits while preserving the consistency of our measure.

A $\sigma$-algebra $\mathcal{F}$ on a set $X$ is a collection of subsets of $X$ such that
$$
X \in \mathcal{F}, \qquad
A \in \mathcal{F} \Rightarrow X \setminus A \in \mathcal{F}, \qquad
A_1,A_2,\dots \in \mathcal{F} \Rightarrow \bigcup_{n=1}^\infty A_n \in \mathcal{F}.
$$
These closure properties ensure that measures are consistent under negation and countable limits.

\paragraph{Example.}
When $X\subset \mathbb{R}^d$, we want
to treat ``geometric'' subsets of $X$ as legitimate events, e.g. open balls, half-spaces,
and threshold sets such as $\{x\in X : f(x)\le a\}$.

We therefore take $\mathcal{F}=\mathcal{B}(X)$, the \textbf{Borel $\sigma$-algebra} on $X$,
defined as the \emph{smallest} $\sigma$-algebra that contains all open subsets of $X$
(equivalently, the $\sigma$-algebra generated by open balls).

Open sets matter here because geometric events (``within radius $r$ of a point'', 
``inside a coverage footprint'') are naturally described using open balls and their unions, 
and we want those events to be measurable by default.

\subsection{Measurable Functions}

\paragraph{Motivation.}
We have considered measuring sets (subsets viewed as events). We now consider measuring functions
that represent physical quantities on $X$ (SINR, received power, interference, etc.), because many events
are defined by thresholding such a function.

\paragraph{Intuition.}
Let $x$ be a random location drawn according to a probability measure $\mu$ on $X$
(for example, $\mu$ could be the distribution of user locations in a latitude band).
If $f:X\to\mathbb{R}$ is a physical quantity (such as SINR), then $f(x)$ is a random real number (a random variable).

Terminology note: in this document, when we say ``distribution'' of $x$ we mean a \emph{probability measure}
$\mu\in\mathcal{P}(X)$, so $\mu(X)=1$. (More general, non-normalised measures also appear in modeling---for example,
Poisson point process \emph{intensity} measures---but probabilities are always defined using a probability measure.)

What we want to quantify are statements like ``subset of locations where the SINR is below the threshold $a$''.
This statement corresponds to the event 
$$
\{x\in X : f(x)\le a\}.
$$
Its probability is the $\mu$-mass of that set:
$$
p = \mu(\{x\in X : f(x)\le a\}).
$$

\paragraph{Definition.}
To ensure this definition is consistent (and to allow integrals like averages), we need the set
$\{x\in X : f(x)\le a\}$ to be an element of $\mathcal{F}$ for every threshold $a$.
This is exactly what it means for $f$ to be measurable.

Formally, we equip the domain with a measurable space $(X,\mathcal{F})$ and the real line with
its Borel $\sigma$-algebra $(\mathbb{R},\mathcal{B}(\mathbb{R}))$.
A function $f:X\to\mathbb{R}$ is measurable if for every Borel set $A\in\mathcal{B}(\mathbb{R})$,
the preimage $f^{-1}(A)\in\mathcal{F}$.
In practice, it is enough to check threshold sets: $f$ is measurable if for every $a\in\mathbb{R}$,
$$
\{x\in X : f(x)\le a\}\in\mathcal{F}.
$$

\paragraph{Example.}
Finally, measurability is what makes expectations meaningful: if $f$ is integrable,
the average value of the physical quantity is $\int_X f(x)\,d\mu(x)$.
When $\mathcal{F}=\mathcal{B}(X)$ and $f$ is continuous (as for most geometric/radio quantities),
$f$ is automatically measurable.

\subsection{Push-Forward Measures}
\paragraph{Motivation.}
The notion of a push-forward measure explains how a measure on $X$ (such as the distribution of locations)
induces a corresponding measure on another space through a function. In our setting this ``other space'' is often
$\mathbb{R}$ (e.g., the real-valued SINR), but the same construction applies to maps $X\to Y$ (as in optimal transport).

\paragraph{Intuition.}
Start with a random location $x\sim \mu$ on $X$. Apply a deterministic map $f$ to obtain a new value $y=f(x)$
(for example a real-valued SINR, or a transported location in $Y$). Even though $f$ is deterministic, the output is
random because the input $x$ is random; the push-forward is simply the distribution of $y$.

\paragraph{Definition.}
Formally, let $(X,\mathcal{F})$ and $(Y,\mathcal{G})$ be measurable spaces, let $\mu$ be a measure on $(X,\mathcal{F})$,
and let $f:X\to Y$ be measurable. The \textbf{push-forward} (also called \textbf{image}) measure $f_{\#}\mu$ on
$(Y,\mathcal{G})$ is defined by
$$
(f_{\#}\mu)(B) := \mu(f^{-1}(B)), \qquad B\in\mathcal{G}.
$$
If $\mu$ is a probability measure, then $f_{\#}\mu$ is again a probability measure.
This says: the mass assigned to a set $B$ of outputs equals the mass of the set of inputs that map into $B$.

\paragraph{Example.}
In the common case $Y=\mathbb{R}$ with $\mathcal{G}=\mathcal{B}(\mathbb{R})$, letting $B=(-\infty,a]$ gives
$$
(f_{\#}\mu)((-\infty,a]) = \mu(\{x\in X : f(x)\le a\}),
$$
which is exactly the probability statement we use for threshold events.

\subsubsection{Push-Forward in the Optimal Transport Problem}
\paragraph{Motivation.}
In optimal transport, we want to \emph{move mass} from a source distribution $\mu$ on a space $X$ (a ``pile'')
to a target distribution $\nu$ on a space $Y$ (a ``hole''), while minimizing a cost (the objective).

\paragraph{Intuition.}
A transport map $T:X\to Y$ tells us where each location $x$ sends its mass. The requirement that ``the pile becomes
the hole'' is exactly that transporting $\mu$ through $T$ yields the target distribution $\nu$.

\paragraph{Definition.}
We say that $T$ \emph{pushes $\mu$ forward to $\nu$} if
$$
T_{\#}\mu=\nu,
$$
meaning that for every measurable set $B\subseteq Y$,
$$
\nu(B)=\mu(T^{-1}(B)).
$$
This is exactly the push-forward definition from the previous subsection, with the specific map $f$ renamed to $T$
(to emphasize that it is a transport map) and the codomain being $Y$ rather than $\mathbb{R}$.

\textbf{Kantorovich formulation.} The Monge formulation above restricts us to a \emph{deterministic} assignment:
each $x$ sends all its mass to a single point $T(x)$. This can be too restrictive (for example, if a point of $\mu$
must supply multiple points in $\nu$) and it leads to a nonconvex optimization problem.

The Kantorovich idea is to relax ``one destination per source point'' into a \emph{transport plan} that can split mass:
instead of choosing a map $T$, we choose a joint measure $\gamma$ on $X\times Y$ whose mass at $(x,y)$ represents
how much mass is transported from $x$ to $y$.

In optimisation language, the \textbf{decision variable} (also called the \textbf{optimisation variable}) is the
mathematical object we are free to choose in order to minimise the objective under constraints.
In the Kantorovich formulation, the decision variable is the joint measure $\gamma$ on $X\times Y$.

To represent a valid transport plan from $\mu$ to $\nu$, $\gamma$ should have marginals $\mu$ and $\nu$.
Equivalently, the constraints that $\gamma$ has marginals equal to $\mu$ and $\nu$ can be written as push-forward relations
$$
(\pi_X)_{\#}\gamma=\mu,\qquad (\pi_Y)_{\#}\gamma=\nu,
$$
where $\pi_X(x,y)=x$ and $\pi_Y(x,y)=y$ are the coordinate projections.
These are exactly push-forwards of the same form $T_{\#}(\cdot)$, with the ``map'' $T$ being instantiated as
$T=\pi_X$ and $T=\pi_Y$, respectively.

To make the origin and endpoint explicit: $\gamma$ is a measure on the \emph{product space} $X\times Y$.
The map $\pi_X:X\times Y\to X$ takes a pair $(x,y)$ and returns its $X$-coordinate $x$, so $(\pi_X)_{\#}\gamma$
is a measure \emph{on $X$}. Likewise, $\pi_Y:X\times Y\to Y$ returns the $Y$-coordinate, so $(\pi_Y)_{\#}\gamma$
is a measure \emph{on $Y$}.

\paragraph{Example.}
Let $X=[0,1]$, $Y=[0,2]$, and let $\mu$ be the uniform distribution on $[0,1]$. Define $T(x)=2x$.
Then $T_{\#}\mu$ is the uniform distribution on $[0,2]$ (the interval is stretched, so density halves), i.e.\ $\nu=T_{\#}\mu$.

Kantorovich-only example (mass splitting). Let $X=\{0\}$ and $Y=\{-1,+1\}$. Let $\mu=\delta_0$ (all mass at $0$) and
let $\nu=\tfrac12\delta_{-1}+\tfrac12\delta_{+1}$. There is no map $T:X\to Y$ such that $T_{\#}\mu=\nu$, because any map
must send all mass to a single point in $Y$. But there is a transport plan $\gamma$ on $X\times Y$ defined by
$$
\gamma(0,-1)=\tfrac12,\qquad \gamma(0,+1)=\tfrac12,
$$
which satisfies $(\pi_X)_{\#}\gamma=\mu$ and $(\pi_Y)_{\#}\gamma=\nu$.

\section{Optimal Transport and Cournot-Nash Equilibria}
This section summarises the background and key ideas in \cite{blanchet2016optimal} in a format aligned with our study.

We are interested in settings with many agents, each attempting to optimise for itself. What sorts of equilibria arise?
What do these equilibria look like? When do they exist, and when are they unique?

For example, consider a large number of users (agents) choosing which satellite (action) to connect to. 
Each user's utility depends on which satellite it chooses, but also on how many other users choose the 
same satellite (congestion externality). We wish to characterise the equilibria of this system, so we may 
understand not only how users behave but also how to design the system (e.g., satellite parameters) to 
achieve desirable outcomes.

Each user is solving an individual optimisation problem, but the problems are coupled through the congestion effects.
We need to characterise jointly the optimisation problem as well as the equilibrium conditions. It would not suffice 
to simply solve each user's optimisation problem in isolation, or to show that an equilibrium exists. It is essential 
to simultaneously characterise individual optimisation strategies and the resulting equilibrium.

In game-theoretic terms, we are interested in \textbf{non-atomic games with externalities}, where each agent is infinitesimally
small (non-atomic) and the cost incurred by each agent depends on the aggregate behaviour of all agents (externalities).

\subsection{Problem Class, Background, and Gap}

\cite{blanchet2016optimal} studies non-atomic games with externalities and gives classes of models where Cournot--Nash
equilibria are not only shown to exist, but can also be characterised in a way that supports uniqueness results and
computation. Their stated gap is that, while abstract existence results for Cournot--Nash equilibria are well known, they
do not typically yield a tractable characterisation (or numerical scheme) for the equilibrium distribution.

\paragraph{Where do externalities and congestion come from?}
In many applications, an agent's cost depends on what \emph{others} do through an aggregate quantity. For instance, in a
wireless association model, if many users choose the same satellite (or the same beam), then service quality degrades;
in road choice, if many drivers choose the same road, travel time increases. These effects are called \emph{externalities}
because an agent's cost depends on the aggregate behaviour of the population, not only on the agent's own action.
The paper models such effects through an action distribution $\nu$ and an externality term $V[\nu](y)$ that modifies the
cost of playing action $y$.

However, once we introduce such externalities (especially congestion-type effects), there is a key technical
obstruction: the dependence of costs on the population distribution can fail to have the regularity properties assumed
by standard equilibrium existence proofs.

\paragraph{What is the ``obstruction''?}
Here, an \emph{obstruction} simply means a technical barrier that prevents us from applying the most common existence
proof techniques ``out of the box''. Concretely, many existence proofs for equilibria rely on a fixed-point argument and
therefore need a basic stability property: if the population action distribution $\nu$ changes a little, then the cost
landscape faced by agents should also change a little. Without such stability, it becomes hard to even define a
well-behaved ``best-response of the population'' map, let alone prove it has a fixed point.

Congestion effects are often \emph{local}: the cost at action $y$ depends on how crowded that action is. A common model is
$V[\nu](y)=f(\nu(y))$, where $\nu(y)$ denotes a density. The difficulty is that densities are fragile: two action
distributions can look almost the same when you only look at coarse averages (e.g.\ averages over regions), yet have very
different pointwise densities because mass can concentrate into narrow spikes or oscillate rapidly. If $V[\nu](y)$ depends
directly on $\nu(y)$, these small ``distribution-level'' changes can produce large changes in the local congestion cost,
so the regularity assumptions behind standard fixed-point existence arguments fail.

To make these statements precise, we first fix the paper's equilibrium notation (Subsection~\ref{sec:cn_equilibrium}) and then explain
the relevant notion of ``two distributions being close'' (weak-$*$ topology) and why local congestion breaks it
(Subsection~\ref{sec:cn_weakstar_gap}).

The key contribution in \cite{blanchet2016optimal} is to exploit an optimal-transport structure (and a variational formulation in potential
cases) to turn equilibrium computation into an optimisation problem over distributions.



\subsection{Equilibrium}\label{sec:cn_equilibrium}
In a finite-player game, one can describe an equilibrium by listing each player's action (or mixed strategy).
In a \emph{non-atomic} game with a continuum of agents, this is not the right object: we do not label agents individually,
and a single agent has negligible influence on the aggregate outcome. Instead, we describe equilibria at the population
level using probability measures.

\subsubsection{Spaces and Probability Measures}
In the paper's notation, $X$ is the \emph{type space} (agent characteristics) and $Y$ is the \emph{action space}.
We write $\mathcal{P}(X)$ for the set of \emph{Borel} probability measures on $X$ and similarly $\mathcal{P}(Y)$.
Here ``Borel'' means the measures are defined on the Borel $\sigma$-algebra generated by the topology on $X$ (and $Y$),
which is the standard choice when $X$ and $Y$ are compact metric spaces. (The ``continuum of agents'' aspect is captured
by describing populations via measures.)
The population of types is given exogenously by a probability measure $\mu\in\mathcal{P}(X)$.

\subsubsection{Equilibrium as a Joint Distribution}
What matters is the population-level assignment of actions to types: for each type $x$, what fraction of agents choose
each action $y$. This assignment is represented by a joint probability measure on $X\times Y$.

An equilibrium is described by a joint probability measure $\gamma\in\mathcal{P}(X\times Y)$, where (informally)
$\gamma(A\times B)$ is the probability that an agent has type in $A\subseteq X$ and chooses an action in $B\subseteq Y$.
The first marginal of $\gamma$ is the induced distribution of types and must equal $\mu$.
The second marginal of $\gamma$ is the induced distribution of actions, denoted $\nu\in\mathcal{P}(Y)$.
Using the coordinate projections $\pi_X(x,y)=x$ and $\pi_Y(x,y)=y$, these marginal constraints can be written as
push-forward relations:
$$
(\pi_X)_{\#}\gamma=\mu,\qquad \nu=(\pi_Y)_{\#}\gamma.
$$
If the equilibrium is \emph{pure}, then there is a map $T:X\to Y$ such that $\gamma$ is carried by the graph of $T$,
i.e.\ $\gamma=(\mathrm{id},T)_{\#}\mu$.

\subsubsection{Costs and Externalities}
Costs are additively separable: an $x$-type agent who plays action $y$ incurs a cost of the form
$$
\Phi(x,y,\nu)=c(x,y)+V[\nu](y),
$$
where $c(x,y)$ is an idiosyncratic term (distance, mismatch, etc.), and the map $\nu\mapsto V[\nu]$ captures
externalities induced by the aggregate action distribution $\nu$.

\subsubsection{Reference Measure, Densities, and Local Congestion}
To speak about ``local density'' at an action $y$, we need $\nu$ to admit a density with respect to some reference
measure on $Y$.
Fix a reference measure $m_0$ on $Y$ and restrict to action distributions $\nu$ that are absolutely continuous with
respect to $m_0$ (written $\nu\ll m_0$). In that case, $\nu$ admits a density $\frac{d\nu}{dm_0}$, and with a slight abuse
of notation we write $\nu(y)$ for this density value at $y$.
Local congestion costs take the form
$$
V[\nu](y)=f(\nu(y)),
$$
so the externality cost at $y$ depends pointwise on how concentrated the action distribution is at $y$.

\subsubsection{Cournot--Nash Equilibrium Condition (Formal Definition)}\label{sec:cn_equilibrium_definition}
\paragraph{Motivation.}
So far, $\gamma$ is just a joint distribution that matches the exogenous type distribution $\mu$ and induces an action
distribution $\nu$. To call $\gamma$ an \emph{equilibrium}, we must encode the behavioural requirement that agents choose
cost-minimising actions \emph{given} the aggregate action distribution $\nu$ that results from everyone's choices.

\paragraph{Intuition.}
Fix a candidate aggregate action distribution $\nu$. Each type $x$ then faces a deterministic optimisation problem:
choose $y\in Y$ to minimise $c(x,y)+V[\nu](y)$. An equilibrium is a self-consistent joint distribution $\gamma$ such that
(i) its second marginal is exactly the $\nu$ used in the optimisation, and (ii) $\gamma$ puts mass only on pairs $(x,y)$
where $y$ is a best response for type $x$.

\paragraph{Definition.}
Let $\gamma\in\mathcal{P}(X\times Y)$ and let $\nu=(\pi_Y)_{\#}\gamma$ be its second marginal.
We say that $\gamma$ is a \textbf{Cournot--Nash equilibrium} (in distributional form) if
\begin{enumerate}[label=(\roman*)]
    \item $(\pi_X)_{\#}\gamma=\mu$ (types are distributed according to $\mu$), and
    \item $\gamma$ is concentrated on best responses given $\nu$, meaning
    $$
    \gamma\bigl(\{(x,y)\in X\times Y:\ c(x,y)+V[\nu](y)=\min_{z\in Y}\bigl(c(x,z)+V[\nu](z)\bigr)\}\bigr)=1.
    $$

    In probabilistic terms: let $(\Omega,\mathcal{F},\mathbb{P})$ be a probability space and let $(X_{\Omega},Y_{\Omega})$
    be an $(X\times Y)$-valued random variable with law $\gamma$ (so $\mathbb{P}[(X_{\Omega},Y_{\Omega})\in A]=\gamma(A)$ for
    every measurable $A\subseteq X\times Y$). Define the event
    $$
    E := \Bigl\{c(X_{\Omega},Y_{\Omega})+V[\nu](Y_{\Omega})=\min_{z\in Y}\bigl(c(X_{\Omega},z)+V[\nu](z)\bigr)\Bigr\}.
    $$
    Then the condition $\gamma(\cdot)=1$ above is exactly the statement $\mathbb{P}(E)=1$: with probability one, the sampled
    action $Y_{\Omega}$ is a best response for the sampled type $X_{\Omega}$ given the aggregate distribution $\nu$.

\end{enumerate}


\paragraph{Intuition.}
This expresses a ``best response for almost every agent'' condition in a non-atomic (continuum) model: we do \emph{not}
quantify over individually labelled agents. Instead, we quantify over the population measure $\gamma$: the set of
non-best-response pairs has $\gamma$-mass $0$, i.e.\ if we pick an agent at random from the population (sample
$(x,y)\sim\gamma$), then with probability $1$ the realised action $y$ is a best response for the realised type $x$.

Equivalently, if we define the best-response set $B_{\nu}\subseteq X\times Y$ by the bracketed condition inside
$\gamma(\cdot)$ above, then $\gamma(B_{\nu})=1$ and the event $E$ is exactly $\{(X_{\Omega},Y_{\Omega})\in B_{\nu}\}$.

\paragraph{Alternative Formulation.}

Equivalently, define the \emph{value function}
$$
\varphi(x):=\min_{z\in Y}\bigl(c(x,z)+V[\nu](z)\bigr),
$$
then the equilibrium condition above can be written as the pair of inequalities
$$
c(x,y)+V[\nu](y)\ge \varphi(x)\quad\text{for all }x\in X\text{ and (where defined) }y\in Y,
$$
with equality holding $\gamma$-almost everywhere (i.e.\ for $\gamma$-almost every pair $(x,y)$).

\paragraph{Example.}
In a satellite-association toy model, let $x$ be a user location and let $y$ index a satellite/beam choice. The term
$c(x,y)$ can represent a geometry-based penalty (e.g.\ path loss), and $V[\nu](y)$ can represent a congestion penalty for
how many users choose satellite/beam $y$ (captured by the aggregate action distribution $\nu$). The equilibrium
condition says that, under the self-consistent distribution of choices, almost every user is assigned only to a choice
that minimises its own cost given the congestion created by all users.


\paragraph{Motivation.}
The equilibrium definition in Subsubsection~\ref{sec:cn_equilibrium_definition} is a self-consistent best-response
condition on a joint distribution $\gamma$.
This definition is conceptually clear, but it does not yet tell us how to \emph{compute} an equilibrium, or how to obtain
uniqueness and stability properties.

\paragraph{Intuition.}
The paper's key step is to identify an important class of models where the equilibrium conditions are exactly the
first-order optimality conditions of a single global objective. In that case, equilibrium computation becomes an
optimisation problem over distributions. Optimal transport enters through the transport cost term, while congestion and
interactions enter through an ``energy'' functional of the action distribution.

\paragraph{Definition.}
In the \emph{potential} (variational) case, the externality map $\nu\mapsto V[\nu]$ is assumed to be the first variation
of some functional $\mathcal{E}$ defined on action distributions:
$$
V[\nu] \;=\; \frac{\delta \mathcal{E}}{\delta \nu}.
$$
Under this assumption, equilibrium action distributions $\nu$ can be characterised as minimisers of an OT-structured
functional of the form
$$
\nu \in \arg\min_{\tilde\nu}\ \mathcal{W}_c(\mu,\tilde\nu) + \mathcal{E}[\tilde\nu],
$$
and an equilibrium coupling $\gamma$ is then obtained as an optimal transport plan between $\mu$ and that minimising
$\nu$.

\paragraph{Example.}
A standard ``good'' potential form used in the paper is
$$
\mathcal{E}[\nu] \;=\; \int_Y F(y,\nu(y))\,dm_0(y) \;+\; \frac12\iint_{Y\times Y} \phi(y,z)\,d\nu(y)\,d\nu(z),
$$
where the first term is a local congestion energy (with $F'(\cdot)=f(\cdot)$) and the second term is a symmetric pairwise
interaction energy (symmetry is what makes it a potential).

\subsection{Why Variational Structure Matters}\label{sec:cn_variational_importance}

\paragraph{Motivation.}
The ability to write equilibrium computation as a global minimisation problem is the reason this approach becomes
practical. Without a variational structure, we typically lose both tractability and the clean uniqueness/stability
statements that come from convexity (or convexity along transport geodesics).

\paragraph{Intuition.}
For general Nash equilibria, there is no reason that ``the equilibrium'' should be the minimiser of a single scalar
objective: different players optimise different objectives, and their incentives can create cycles and rotational
behaviour. A variational (potential) structure is exactly the condition that makes the equilibrium conditions integrable
into a single energy landscape: equilibria are stationary points of that landscape, and in the convex cases they are
global minimisers.

In the Cournot--Nash setting considered by \cite{blanchet2016optimal}, the variational structure is what turns a
definition-by-best-response (Subsubsection~\ref{sec:cn_equilibrium_definition}) into a computable characterisation
(Subsection~\ref{sec:cn_variational_formulation}): solve a minimisation problem for the action distribution $\nu$, then
recover the coupling $\gamma$ using optimal transport.

\paragraph{Definition.}
We will refer to the model as \textbf{variational} (or \textbf{potential}) if there exists a functional
$\mathcal{E}$ on action distributions such that $V[\nu]=\delta\mathcal{E}/\delta\nu$.
In this case, the equilibrium conditions coincide with the first-order optimality conditions of the global objective
$\mathcal{W}_c(\mu,\nu)+\mathcal{E}[\nu]$.
If, in addition, this objective is (geodesically) convex in a suitable sense, then the equilibrium is unique and stable
to perturbations of the inputs $(\mu,c,\mathcal{E})$.

\paragraph{Example.}
Not every game has this structure. In finite games such as Rock--Paper--Scissors or Matching Pennies, incentives create
cyclic dominance, and there is no single potential function whose value strictly improves along unilateral best-response
moves. In continuous-action games, strong asymmetric cross-effects can lead to non-monotone best-response operators (the
system behaves more like a rotation than a gradient flow). The relevance here is conceptual: the Blanchet--Carlier
approach relies on being able to convert equilibrium computation into minimisation, which is precisely what fails when a
game is not variational.

\subsection{Practical Takeaway: Variational Characterisation vs.\ Continuum Complexity}
\label{sec:cn_weakstar_gap}
\label{sec:cn_variational_formulation}
The main technical payoff in \cite{blanchet2016optimal} is that, in the \emph{potential} (variational) case, equilibrium
computation can be reframed as a single global optimisation problem (an optimal-transport assignment term plus an energy
term in the action distribution). This provides a principled \emph{characterisation} of equilibrium and, in favourable
convex cases, supports uniqueness and stable numerical schemes.

For this study, however, pursuing the full continuum route is not an attractive path: the formulation lives on
continuum objects (measures on compact metric spaces, densities, and associated regularity/continuity assumptions), while
our satellite questions are naturally finite and simulator-driven. We therefore retain the equilibrium definition as a
useful theoretical lens, but we do not attempt to implement or rely on the full continuum variational machinery.



% References for citations used here are provided by the parent document (`docs/spec.tex`).
`.
% It intentionally contains no `\documentclass`, no preamble, and no `\begin{document}`.

\section{Scope}
We treat \textbf{planning} as the strategic placement and configuration of available
infrastructure under constraints such as interference, and \textbf{dimensioning} as the
minimum resource sizing required to satisfy per-user QoS targets under stochastic
load and resource outages.
Our three core problems of interest are interference, per-user QoS/outage, and route existence.

\section{Common Model Assumptions}
\begin{itemize}
    \item \textbf{Users} form a spatial point process on the Earth surface, within
    a latitude band.
    \item \textbf{Satellites} form a spatial point process on an orbital shell,
    gated by visibility.
    \item \textbf{Radio model} uses a received-power formulation
    $P(x,y)=P_t G_{xy}\|x-y\|^{-\gamma}$ with Rayleigh fading
    $G_{xy}\sim \mathrm{Exp}(1)$ and SINR-based capacity $C=\mathcal{W}\log_2(1+\mathrm{SINR})$.
    This follows standard stochastic-geometry modeling conventions \cite{gomez2018stochastic}.
\end{itemize}

Definitions: $P_t$ is per-satellite transmit power (W); $G_{xy}$ is a unitless
small-scale fading power gain for the link $x\leftrightarrow y$; $\|x-y\|$ is
3D link distance (m); $\gamma$ is the path-loss exponent; $\mathcal{W}$ is bandwidth (Hz);
$\mathrm{SINR}=S/(I+N)$ with $S$ the received signal power, $I$ aggregate
interference, and $N=N_0\mathcal{W}$ the noise power with $N_0$ noise spectral density (W/Hz).

\section{Planning: Interference as the Primary Objective}
We view planning as choosing configuration variables (e.g., satellite density,
altitude, power, bandwidth, and visibility constraints) so that variables of interest are
controlled (primarily intereference) while maintaining acceptable service quality.
\begin{itemize}
    \item \textbf{Interference metric}: for user at $x$ with serving satellite
    $y_0$, define $I(x)=\sum_{y\in \Phi_b\setminus\{y_0\}} P(x,y)$, that is, the aggregate power
    received from all visible satellites except the serving one.
    \item \textbf{Planning target}: enforce a region-level constraint such as
    $\mathbb{P}(\mathrm{SIR}(x)\ge\beta)\ge p_{\mathrm{sir}}$ or
    $\mathbb{E}[I(x)]\le I_{\max}$ for users in the band.
    \item \textbf{Equilibrium framing}: treat operators or satellites as agents
    whose actions affect aggregate interference, then study equilibria under
    congestion penalties.
\end{itemize}

\section{Dimensioning: Per-User QoS Under Network Outage}
Dimensioning seeks the minimum satellite intensity $N$ such that per-user QoS
targets are met with high confidence, even when the network is resource-limited.
\begin{itemize}
    \item \textbf{Per-user outage}: $O_i=1$ if a user is unserved or if
    throughput $C_i < R_{\min}$; otherwise $O_i=0$.
    \item \textbf{Outage rate}: $o = \frac{1}{|\mathcal{U}|}\sum_{i\in\mathcal{U}} O_i$.
    \item \textbf{Dimensioning target}: find the smallest $N$ such that
    $\mathbb{P}(o \le o_{\max}) \ge p_{\mathrm{target}}$ under the stated priors.
\end{itemize}

\section{Routing: Route Existence Under Time-Varying Connectivity}
We study whether an end-to-end route exists between endpoints given a time-varying visibility/connectivity graph
(ground--satellite and inter-satellite links), reporting $\mathbb{P}(\text{route exists})$ (or a certified lower bound).

\section{Planned Experiments}
\begin{itemize}
    \item Vary satellite intensity and compute interference distributions for
    latitude bands to identify feasible planning regimes.
    \item Measure per-user outage and certify feasibility over Monte Carlo trials,
    reporting the minimum $N$ that meets $(o_{\max}, p_{\mathrm{target}})$.
    \item Measure route existence (as a graph-connectivity event) under representative visibility/connectivity snapshots.
\end{itemize}



\section{Theoretical Structure}

\subsection{Topological and Metric Spaces}

\paragraph{Motivation.}
Many sets we work with in applications (such as subsets of $\mathbb{R}^d$) come with a natural notion of
``closeness'' and ``continuity''. This structure is what lets us talk about open sets, continuous costs, and Borel sets,
which in turn are the standard measurable sets used for probability measures.

\paragraph{Intuition.}
A \emph{topology} tells us which subsets of a set should be considered ``open'', generalising open intervals in
$\mathbb{R}$ and open balls in $\mathbb{R}^d$. A \emph{metric} is a concrete way to specify closeness by giving distances.
Compactness is a finiteness condition (no ``escape to infinity'') that makes many existence arguments behave well.

\paragraph{Definition.}
A \textbf{topological space} is a pair $(X,\tau)$ where $\tau\subseteq 2^X$ is a collection of subsets of $X$ (called
\emph{open sets}) such that:
$$
\emptyset\in\tau,\quad X\in\tau,\quad
\bigcup_{i\in I} U_i \in\tau\ \text{for any family }(U_i)_{i\in I}\subseteq\tau,\quad
U\cap V\in\tau\ \text{for all }U,V\in\tau.
$$

A \textbf{metric space} is a pair $(X,d)$ where $d:X\times X\to[0,\infty)$ satisfies (i) $d(x,y)=0\Leftrightarrow x=y$,
(ii) symmetry, and (iii) the triangle inequality. Every metric $d$ induces a topology by taking open sets to be unions of
open balls $B_r(x)=\{z\in X:d(z,x)<r\}$.

A metric space $(X,d)$ is \textbf{compact} if every open cover of $X$ has a finite subcover (equivalently in metric
spaces, every sequence has a convergent subsequence). A \textbf{compact metric space} is a metric space that is compact.

\paragraph{Example.}
$X=[0,1]\subset\mathbb{R}$ with the usual distance $d(x,y)=|x-y|$ is a compact metric space.
More generally, every closed and bounded subset of $\mathbb{R}^d$ with the Euclidean distance is compact.

\subsection{Measurability}

\paragraph{Motivation.}
We wish to make quantitative statements (such as probabilities of events). In order to do so, we
must ensure that the sets we are making statements about are measurable, that is, we must ensure that 
a simple real number can be assigned to their subsets in a consistent manner that correctly reflects their 
relative sizes (masses). For example, we wish to talk about probabilities of events such as 
``signal to noise and interference is below threshold with probability''.

\paragraph{Intuition.}
How do we create a framework that allows us to measure subsets of a set (e.g. events in the case of probability)? 

\paragraph{Definition.}
A \textbf{measurable space} is a pair $(X,\mathcal{F})$ where $X$ is a set and $\mathcal{F}$ is a collection of subsets
of $X$ (called ``events'') on which we will be allowed to assign sizes.
A \textbf{measure} on $(X,\mathcal{F})$ is a function $\mu:\mathcal{F}\to[0,\infty]$ such that $\mu(\emptyset)=0$ and
for any pairwise disjoint sets $A_1,A_2,\dots \in \mathcal{F}$,
$$
\mu\!\left(\bigcup_{n=1}^\infty A_n\right)=\sum_{n=1}^\infty \mu(A_n).
$$
A \textbf{probability measure} is a measure with $\mu(X)=1$.

\paragraph{Example.}
If $X=\{1,2,\dots,n\}$ is finite, a natural choice is $\mathcal{F}=2^X$ (all subsets), and the uniform probability
measure $\mu(A)=|A|/n$. Then $\mu(A)$ is the probability of the event ``the outcome lies in $A$''.

\subsection{$\sigma$-algebras}

\paragraph{Motivation.}
In probability, we want to manipulate events using the usual set operations: we negate events, take unions of events,
and take limits of event sequences. The collection of ``allowed'' events should therefore be closed under these
operations, so that we do not leave the space of events when doing algebra.

\paragraph{Intuition.}
We do \emph{not} assume every subset of $X$ is measurable: we pick a structured family of subsets that is large enough
for the geometric/physical events we care about, but stable under complements and countable unions.

\paragraph{Definition.}
This is done by equipping these sets (e.g. $X$ the set of all agent types, $Y$ the set of all actions) with a 
$\sigma$-algebra $\mathcal{F}$, which is a collection of subsets defined in such a way that assigning a real number 
to each subset in $\mathcal{F}$ is consistent with the usual rules of set operations (union, intersection, complement) 
and countable limits. In other words, we may add, subtract, and take limits while preserving the consistency of our measure.

A $\sigma$-algebra $\mathcal{F}$ on a set $X$ is a collection of subsets of $X$ such that
$$
X \in \mathcal{F}, \qquad
A \in \mathcal{F} \Rightarrow X \setminus A \in \mathcal{F}, \qquad
A_1,A_2,\dots \in \mathcal{F} \Rightarrow \bigcup_{n=1}^\infty A_n \in \mathcal{F}.
$$
These closure properties ensure that measures are consistent under negation and countable limits.

\paragraph{Example.}
When $X\subset \mathbb{R}^d$, we want
to treat ``geometric'' subsets of $X$ as legitimate events, e.g. open balls, half-spaces,
and threshold sets such as $\{x\in X : f(x)\le a\}$.

We therefore take $\mathcal{F}=\mathcal{B}(X)$, the \textbf{Borel $\sigma$-algebra} on $X$,
defined as the \emph{smallest} $\sigma$-algebra that contains all open subsets of $X$
(equivalently, the $\sigma$-algebra generated by open balls).

Open sets matter here because geometric events (``within radius $r$ of a point'', 
``inside a coverage footprint'') are naturally described using open balls and their unions, 
and we want those events to be measurable by default.

\subsection{Measurable Functions}

\paragraph{Motivation.}
We have considered measuring sets (subsets viewed as events). We now consider measuring functions
that represent physical quantities on $X$ (SINR, received power, interference, etc.), because many events
are defined by thresholding such a function.

\paragraph{Intuition.}
Let $x$ be a random location drawn according to a probability measure $\mu$ on $X$
(for example, $\mu$ could be the distribution of user locations in a latitude band).
If $f:X\to\mathbb{R}$ is a physical quantity (such as SINR), then $f(x)$ is a random real number (a random variable).

Terminology note: in this document, when we say ``distribution'' of $x$ we mean a \emph{probability measure}
$\mu\in\mathcal{P}(X)$, so $\mu(X)=1$. (More general, non-normalised measures also appear in modeling---for example,
Poisson point process \emph{intensity} measures---but probabilities are always defined using a probability measure.)

What we want to quantify are statements like ``subset of locations where the SINR is below the threshold $a$''.
This statement corresponds to the event 
$$
\{x\in X : f(x)\le a\}.
$$
Its probability is the $\mu$-mass of that set:
$$
p = \mu(\{x\in X : f(x)\le a\}).
$$

\paragraph{Definition.}
To ensure this definition is consistent (and to allow integrals like averages), we need the set
$\{x\in X : f(x)\le a\}$ to be an element of $\mathcal{F}$ for every threshold $a$.
This is exactly what it means for $f$ to be measurable.

Formally, we equip the domain with a measurable space $(X,\mathcal{F})$ and the real line with
its Borel $\sigma$-algebra $(\mathbb{R},\mathcal{B}(\mathbb{R}))$.
A function $f:X\to\mathbb{R}$ is measurable if for every Borel set $A\in\mathcal{B}(\mathbb{R})$,
the preimage $f^{-1}(A)\in\mathcal{F}$.
In practice, it is enough to check threshold sets: $f$ is measurable if for every $a\in\mathbb{R}$,
$$
\{x\in X : f(x)\le a\}\in\mathcal{F}.
$$

\paragraph{Example.}
Finally, measurability is what makes expectations meaningful: if $f$ is integrable,
the average value of the physical quantity is $\int_X f(x)\,d\mu(x)$.
When $\mathcal{F}=\mathcal{B}(X)$ and $f$ is continuous (as for most geometric/radio quantities),
$f$ is automatically measurable.

\subsection{Push-Forward Measures}
\paragraph{Motivation.}
The notion of a push-forward measure explains how a measure on $X$ (such as the distribution of locations)
induces a corresponding measure on another space through a function. In our setting this ``other space'' is often
$\mathbb{R}$ (e.g., the real-valued SINR), but the same construction applies to maps $X\to Y$ (as in optimal transport).

\paragraph{Intuition.}
Start with a random location $x\sim \mu$ on $X$. Apply a deterministic map $f$ to obtain a new value $y=f(x)$
(for example a real-valued SINR, or a transported location in $Y$). Even though $f$ is deterministic, the output is
random because the input $x$ is random; the push-forward is simply the distribution of $y$.

\paragraph{Definition.}
Formally, let $(X,\mathcal{F})$ and $(Y,\mathcal{G})$ be measurable spaces, let $\mu$ be a measure on $(X,\mathcal{F})$,
and let $f:X\to Y$ be measurable. The \textbf{push-forward} (also called \textbf{image}) measure $f_{\#}\mu$ on
$(Y,\mathcal{G})$ is defined by
$$
(f_{\#}\mu)(B) := \mu(f^{-1}(B)), \qquad B\in\mathcal{G}.
$$
If $\mu$ is a probability measure, then $f_{\#}\mu$ is again a probability measure.
This says: the mass assigned to a set $B$ of outputs equals the mass of the set of inputs that map into $B$.

\paragraph{Example.}
In the common case $Y=\mathbb{R}$ with $\mathcal{G}=\mathcal{B}(\mathbb{R})$, letting $B=(-\infty,a]$ gives
$$
(f_{\#}\mu)((-\infty,a]) = \mu(\{x\in X : f(x)\le a\}),
$$
which is exactly the probability statement we use for threshold events.

\subsubsection{Push-Forward in the Optimal Transport Problem}
\paragraph{Motivation.}
In optimal transport, we want to \emph{move mass} from a source distribution $\mu$ on a space $X$ (a ``pile'')
to a target distribution $\nu$ on a space $Y$ (a ``hole''), while minimizing a cost (the objective).

\paragraph{Intuition.}
A transport map $T:X\to Y$ tells us where each location $x$ sends its mass. The requirement that ``the pile becomes
the hole'' is exactly that transporting $\mu$ through $T$ yields the target distribution $\nu$.

\paragraph{Definition.}
We say that $T$ \emph{pushes $\mu$ forward to $\nu$} if
$$
T_{\#}\mu=\nu,
$$
meaning that for every measurable set $B\subseteq Y$,
$$
\nu(B)=\mu(T^{-1}(B)).
$$
This is exactly the push-forward definition from the previous subsection, with the specific map $f$ renamed to $T$
(to emphasize that it is a transport map) and the codomain being $Y$ rather than $\mathbb{R}$.

\textbf{Kantorovich formulation.} The Monge formulation above restricts us to a \emph{deterministic} assignment:
each $x$ sends all its mass to a single point $T(x)$. This can be too restrictive (for example, if a point of $\mu$
must supply multiple points in $\nu$) and it leads to a nonconvex optimization problem.

The Kantorovich idea is to relax ``one destination per source point'' into a \emph{transport plan} that can split mass:
instead of choosing a map $T$, we choose a joint measure $\gamma$ on $X\times Y$ whose mass at $(x,y)$ represents
how much mass is transported from $x$ to $y$.

In optimisation language, the \textbf{decision variable} (also called the \textbf{optimisation variable}) is the
mathematical object we are free to choose in order to minimise the objective under constraints.
In the Kantorovich formulation, the decision variable is the joint measure $\gamma$ on $X\times Y$.

To represent a valid transport plan from $\mu$ to $\nu$, $\gamma$ should have marginals $\mu$ and $\nu$.
Equivalently, the constraints that $\gamma$ has marginals equal to $\mu$ and $\nu$ can be written as push-forward relations
$$
(\pi_X)_{\#}\gamma=\mu,\qquad (\pi_Y)_{\#}\gamma=\nu,
$$
where $\pi_X(x,y)=x$ and $\pi_Y(x,y)=y$ are the coordinate projections.
These are exactly push-forwards of the same form $T_{\#}(\cdot)$, with the ``map'' $T$ being instantiated as
$T=\pi_X$ and $T=\pi_Y$, respectively.

To make the origin and endpoint explicit: $\gamma$ is a measure on the \emph{product space} $X\times Y$.
The map $\pi_X:X\times Y\to X$ takes a pair $(x,y)$ and returns its $X$-coordinate $x$, so $(\pi_X)_{\#}\gamma$
is a measure \emph{on $X$}. Likewise, $\pi_Y:X\times Y\to Y$ returns the $Y$-coordinate, so $(\pi_Y)_{\#}\gamma$
is a measure \emph{on $Y$}.

\paragraph{Example.}
Let $X=[0,1]$, $Y=[0,2]$, and let $\mu$ be the uniform distribution on $[0,1]$. Define $T(x)=2x$.
Then $T_{\#}\mu$ is the uniform distribution on $[0,2]$ (the interval is stretched, so density halves), i.e.\ $\nu=T_{\#}\mu$.

Kantorovich-only example (mass splitting). Let $X=\{0\}$ and $Y=\{-1,+1\}$. Let $\mu=\delta_0$ (all mass at $0$) and
let $\nu=\tfrac12\delta_{-1}+\tfrac12\delta_{+1}$. There is no map $T:X\to Y$ such that $T_{\#}\mu=\nu$, because any map
must send all mass to a single point in $Y$. But there is a transport plan $\gamma$ on $X\times Y$ defined by
$$
\gamma(0,-1)=\tfrac12,\qquad \gamma(0,+1)=\tfrac12,
$$
which satisfies $(\pi_X)_{\#}\gamma=\mu$ and $(\pi_Y)_{\#}\gamma=\nu$.

\section{Optimal Transport and Cournot-Nash Equilibria}
This section summarises the background and key ideas in \cite{blanchet2016optimal} in a format aligned with our study.

We are interested in settings with many agents, each attempting to optimise for itself. What sorts of equilibria arise?
What do these equilibria look like? When do they exist, and when are they unique?

For example, consider a large number of users (agents) choosing which satellite (action) to connect to. 
Each user's utility depends on which satellite it chooses, but also on how many other users choose the 
same satellite (congestion externality). We wish to characterise the equilibria of this system, so we may 
understand not only how users behave but also how to design the system (e.g., satellite parameters) to 
achieve desirable outcomes.

Each user is solving an individual optimisation problem, but the problems are coupled through the congestion effects.
We need to characterise jointly the optimisation problem as well as the equilibrium conditions. It would not suffice 
to simply solve each user's optimisation problem in isolation, or to show that an equilibrium exists. It is essential 
to simultaneously characterise individual optimisation strategies and the resulting equilibrium.

In game-theoretic terms, we are interested in \textbf{non-atomic games with externalities}, where each agent is infinitesimally
small (non-atomic) and the cost incurred by each agent depends on the aggregate behaviour of all agents (externalities).

\subsection{Problem Class, Background, and Gap}

\cite{blanchet2016optimal} studies non-atomic games with externalities and gives classes of models where Cournot--Nash
equilibria are not only shown to exist, but can also be characterised in a way that supports uniqueness results and
computation. Their stated gap is that, while abstract existence results for Cournot--Nash equilibria are well known, they
do not typically yield a tractable characterisation (or numerical scheme) for the equilibrium distribution.

\paragraph{Where do externalities and congestion come from?}
In many applications, an agent's cost depends on what \emph{others} do through an aggregate quantity. For instance, in a
wireless association model, if many users choose the same satellite (or the same beam), then service quality degrades;
in road choice, if many drivers choose the same road, travel time increases. These effects are called \emph{externalities}
because an agent's cost depends on the aggregate behaviour of the population, not only on the agent's own action.
The paper models such effects through an action distribution $\nu$ and an externality term $V[\nu](y)$ that modifies the
cost of playing action $y$.

However, once we introduce such externalities (especially congestion-type effects), there is a key technical
obstruction: the dependence of costs on the population distribution can fail to have the regularity properties assumed
by standard equilibrium existence proofs.

\paragraph{What is the ``obstruction''?}
Here, an \emph{obstruction} simply means a technical barrier that prevents us from applying the most common existence
proof techniques ``out of the box''. Concretely, many existence proofs for equilibria rely on a fixed-point argument and
therefore need a basic stability property: if the population action distribution $\nu$ changes a little, then the cost
landscape faced by agents should also change a little. Without such stability, it becomes hard to even define a
well-behaved ``best-response of the population'' map, let alone prove it has a fixed point.

Congestion effects are often \emph{local}: the cost at action $y$ depends on how crowded that action is. A common model is
$V[\nu](y)=f(\nu(y))$, where $\nu(y)$ denotes a density. The difficulty is that densities are fragile: two action
distributions can look almost the same when you only look at coarse averages (e.g.\ averages over regions), yet have very
different pointwise densities because mass can concentrate into narrow spikes or oscillate rapidly. If $V[\nu](y)$ depends
directly on $\nu(y)$, these small ``distribution-level'' changes can produce large changes in the local congestion cost,
so the regularity assumptions behind standard fixed-point existence arguments fail.

To make these statements precise, we first fix the paper's equilibrium notation (Subsection~\ref{sec:cn_equilibrium}) and then explain
the relevant notion of ``two distributions being close'' (weak-$*$ topology) and why local congestion breaks it
(Subsection~\ref{sec:cn_weakstar_gap}).

The key contribution in \cite{blanchet2016optimal} is to exploit an optimal-transport structure (and a variational formulation in potential
cases) to turn equilibrium computation into an optimisation problem over distributions.



\subsection{Equilibrium}\label{sec:cn_equilibrium}
In a finite-player game, one can describe an equilibrium by listing each player's action (or mixed strategy).
In a \emph{non-atomic} game with a continuum of agents, this is not the right object: we do not label agents individually,
and a single agent has negligible influence on the aggregate outcome. Instead, we describe equilibria at the population
level using probability measures.

\subsubsection{Spaces and Probability Measures}
In the paper's notation, $X$ is the \emph{type space} (agent characteristics) and $Y$ is the \emph{action space}.
We write $\mathcal{P}(X)$ for the set of \emph{Borel} probability measures on $X$ and similarly $\mathcal{P}(Y)$.
Here ``Borel'' means the measures are defined on the Borel $\sigma$-algebra generated by the topology on $X$ (and $Y$),
which is the standard choice when $X$ and $Y$ are compact metric spaces. (The ``continuum of agents'' aspect is captured
by describing populations via measures.)
The population of types is given exogenously by a probability measure $\mu\in\mathcal{P}(X)$.

\subsubsection{Equilibrium as a Joint Distribution}
What matters is the population-level assignment of actions to types: for each type $x$, what fraction of agents choose
each action $y$. This assignment is represented by a joint probability measure on $X\times Y$.

An equilibrium is described by a joint probability measure $\gamma\in\mathcal{P}(X\times Y)$, where (informally)
$\gamma(A\times B)$ is the probability that an agent has type in $A\subseteq X$ and chooses an action in $B\subseteq Y$.
The first marginal of $\gamma$ is the induced distribution of types and must equal $\mu$.
The second marginal of $\gamma$ is the induced distribution of actions, denoted $\nu\in\mathcal{P}(Y)$.
Using the coordinate projections $\pi_X(x,y)=x$ and $\pi_Y(x,y)=y$, these marginal constraints can be written as
push-forward relations:
$$
(\pi_X)_{\#}\gamma=\mu,\qquad \nu=(\pi_Y)_{\#}\gamma.
$$
If the equilibrium is \emph{pure}, then there is a map $T:X\to Y$ such that $\gamma$ is carried by the graph of $T$,
i.e.\ $\gamma=(\mathrm{id},T)_{\#}\mu$.

\subsubsection{Costs and Externalities}
Costs are additively separable: an $x$-type agent who plays action $y$ incurs a cost of the form
$$
\Phi(x,y,\nu)=c(x,y)+V[\nu](y),
$$
where $c(x,y)$ is an idiosyncratic term (distance, mismatch, etc.), and the map $\nu\mapsto V[\nu]$ captures
externalities induced by the aggregate action distribution $\nu$.

\subsubsection{Reference Measure, Densities, and Local Congestion}
To speak about ``local density'' at an action $y$, we need $\nu$ to admit a density with respect to some reference
measure on $Y$.
Fix a reference measure $m_0$ on $Y$ and restrict to action distributions $\nu$ that are absolutely continuous with
respect to $m_0$ (written $\nu\ll m_0$). In that case, $\nu$ admits a density $\frac{d\nu}{dm_0}$, and with a slight abuse
of notation we write $\nu(y)$ for this density value at $y$.
Local congestion costs take the form
$$
V[\nu](y)=f(\nu(y)),
$$
so the externality cost at $y$ depends pointwise on how concentrated the action distribution is at $y$.

\subsubsection{Cournot--Nash Equilibrium Condition (Formal Definition)}\label{sec:cn_equilibrium_definition}
\paragraph{Motivation.}
So far, $\gamma$ is just a joint distribution that matches the exogenous type distribution $\mu$ and induces an action
distribution $\nu$. To call $\gamma$ an \emph{equilibrium}, we must encode the behavioural requirement that agents choose
cost-minimising actions \emph{given} the aggregate action distribution $\nu$ that results from everyone's choices.

\paragraph{Intuition.}
Fix a candidate aggregate action distribution $\nu$. Each type $x$ then faces a deterministic optimisation problem:
choose $y\in Y$ to minimise $c(x,y)+V[\nu](y)$. An equilibrium is a self-consistent joint distribution $\gamma$ such that
(i) its second marginal is exactly the $\nu$ used in the optimisation, and (ii) $\gamma$ puts mass only on pairs $(x,y)$
where $y$ is a best response for type $x$.

\paragraph{Definition.}
Let $\gamma\in\mathcal{P}(X\times Y)$ and let $\nu=(\pi_Y)_{\#}\gamma$ be its second marginal.
We say that $\gamma$ is a \textbf{Cournot--Nash equilibrium} (in distributional form) if
\begin{enumerate}[label=(\roman*)]
    \item $(\pi_X)_{\#}\gamma=\mu$ (types are distributed according to $\mu$), and
    \item $\gamma$ is concentrated on best responses given $\nu$, meaning
    $$
    \gamma\bigl(\{(x,y)\in X\times Y:\ c(x,y)+V[\nu](y)=\min_{z\in Y}\bigl(c(x,z)+V[\nu](z)\bigr)\}\bigr)=1.
    $$

    In probabilistic terms: let $(\Omega,\mathcal{F},\mathbb{P})$ be a probability space and let $(X_{\Omega},Y_{\Omega})$
    be an $(X\times Y)$-valued random variable with law $\gamma$ (so $\mathbb{P}[(X_{\Omega},Y_{\Omega})\in A]=\gamma(A)$ for
    every measurable $A\subseteq X\times Y$). Define the event
    $$
    E := \Bigl\{c(X_{\Omega},Y_{\Omega})+V[\nu](Y_{\Omega})=\min_{z\in Y}\bigl(c(X_{\Omega},z)+V[\nu](z)\bigr)\Bigr\}.
    $$
    Then the condition $\gamma(\cdot)=1$ above is exactly the statement $\mathbb{P}(E)=1$: with probability one, the sampled
    action $Y_{\Omega}$ is a best response for the sampled type $X_{\Omega}$ given the aggregate distribution $\nu$.

\end{enumerate}


\paragraph{Intuition.}
This expresses a ``best response for almost every agent'' condition in a non-atomic (continuum) model: we do \emph{not}
quantify over individually labelled agents. Instead, we quantify over the population measure $\gamma$: the set of
non-best-response pairs has $\gamma$-mass $0$, i.e.\ if we pick an agent at random from the population (sample
$(x,y)\sim\gamma$), then with probability $1$ the realised action $y$ is a best response for the realised type $x$.

Equivalently, if we define the best-response set $B_{\nu}\subseteq X\times Y$ by the bracketed condition inside
$\gamma(\cdot)$ above, then $\gamma(B_{\nu})=1$ and the event $E$ is exactly $\{(X_{\Omega},Y_{\Omega})\in B_{\nu}\}$.

\paragraph{Alternative Formulation.}

Equivalently, define the \emph{value function}
$$
\varphi(x):=\min_{z\in Y}\bigl(c(x,z)+V[\nu](z)\bigr),
$$
then the equilibrium condition above can be written as the pair of inequalities
$$
c(x,y)+V[\nu](y)\ge \varphi(x)\quad\text{for all }x\in X\text{ and (where defined) }y\in Y,
$$
with equality holding $\gamma$-almost everywhere (i.e.\ for $\gamma$-almost every pair $(x,y)$).

\paragraph{Example.}
In a satellite-association toy model, let $x$ be a user location and let $y$ index a satellite/beam choice. The term
$c(x,y)$ can represent a geometry-based penalty (e.g.\ path loss), and $V[\nu](y)$ can represent a congestion penalty for
how many users choose satellite/beam $y$ (captured by the aggregate action distribution $\nu$). The equilibrium
condition says that, under the self-consistent distribution of choices, almost every user is assigned only to a choice
that minimises its own cost given the congestion created by all users.


\paragraph{Motivation.}
The equilibrium definition in Subsubsection~\ref{sec:cn_equilibrium_definition} is a self-consistent best-response
condition on a joint distribution $\gamma$.
This definition is conceptually clear, but it does not yet tell us how to \emph{compute} an equilibrium, or how to obtain
uniqueness and stability properties.

\paragraph{Intuition.}
The paper's key step is to identify an important class of models where the equilibrium conditions are exactly the
first-order optimality conditions of a single global objective. In that case, equilibrium computation becomes an
optimisation problem over distributions. Optimal transport enters through the transport cost term, while congestion and
interactions enter through an ``energy'' functional of the action distribution.

\paragraph{Definition.}
In the \emph{potential} (variational) case, the externality map $\nu\mapsto V[\nu]$ is assumed to be the first variation
of some functional $\mathcal{E}$ defined on action distributions:
$$
V[\nu] \;=\; \frac{\delta \mathcal{E}}{\delta \nu}.
$$
Under this assumption, equilibrium action distributions $\nu$ can be characterised as minimisers of an OT-structured
functional of the form
$$
\nu \in \arg\min_{\tilde\nu}\ \mathcal{W}_c(\mu,\tilde\nu) + \mathcal{E}[\tilde\nu],
$$
and an equilibrium coupling $\gamma$ is then obtained as an optimal transport plan between $\mu$ and that minimising
$\nu$.

\paragraph{Example.}
A standard ``good'' potential form used in the paper is
$$
\mathcal{E}[\nu] \;=\; \int_Y F(y,\nu(y))\,dm_0(y) \;+\; \frac12\iint_{Y\times Y} \phi(y,z)\,d\nu(y)\,d\nu(z),
$$
where the first term is a local congestion energy (with $F'(\cdot)=f(\cdot)$) and the second term is a symmetric pairwise
interaction energy (symmetry is what makes it a potential).

\subsection{Why Variational Structure Matters}\label{sec:cn_variational_importance}

\paragraph{Motivation.}
The ability to write equilibrium computation as a global minimisation problem is the reason this approach becomes
practical. Without a variational structure, we typically lose both tractability and the clean uniqueness/stability
statements that come from convexity (or convexity along transport geodesics).

\paragraph{Intuition.}
For general Nash equilibria, there is no reason that ``the equilibrium'' should be the minimiser of a single scalar
objective: different players optimise different objectives, and their incentives can create cycles and rotational
behaviour. A variational (potential) structure is exactly the condition that makes the equilibrium conditions integrable
into a single energy landscape: equilibria are stationary points of that landscape, and in the convex cases they are
global minimisers.

In the Cournot--Nash setting considered by \cite{blanchet2016optimal}, the variational structure is what turns a
definition-by-best-response (Subsubsection~\ref{sec:cn_equilibrium_definition}) into a computable characterisation
(Subsection~\ref{sec:cn_variational_formulation}): solve a minimisation problem for the action distribution $\nu$, then
recover the coupling $\gamma$ using optimal transport.

\paragraph{Definition.}
We will refer to the model as \textbf{variational} (or \textbf{potential}) if there exists a functional
$\mathcal{E}$ on action distributions such that $V[\nu]=\delta\mathcal{E}/\delta\nu$.
In this case, the equilibrium conditions coincide with the first-order optimality conditions of the global objective
$\mathcal{W}_c(\mu,\nu)+\mathcal{E}[\nu]$.
If, in addition, this objective is (geodesically) convex in a suitable sense, then the equilibrium is unique and stable
to perturbations of the inputs $(\mu,c,\mathcal{E})$.

\paragraph{Example.}
Not every game has this structure. In finite games such as Rock--Paper--Scissors or Matching Pennies, incentives create
cyclic dominance, and there is no single potential function whose value strictly improves along unilateral best-response
moves. In continuous-action games, strong asymmetric cross-effects can lead to non-monotone best-response operators (the
system behaves more like a rotation than a gradient flow). The relevance here is conceptual: the Blanchet--Carlier
approach relies on being able to convert equilibrium computation into minimisation, which is precisely what fails when a
game is not variational.

\subsection{Practical Takeaway: Variational Characterisation vs.\ Continuum Complexity}
\label{sec:cn_weakstar_gap}
\label{sec:cn_variational_formulation}
The main technical payoff in \cite{blanchet2016optimal} is that, in the \emph{potential} (variational) case, equilibrium
computation can be reframed as a single global optimisation problem (an optimal-transport assignment term plus an energy
term in the action distribution). This provides a principled \emph{characterisation} of equilibrium and, in favourable
convex cases, supports uniqueness and stable numerical schemes.

For this study, however, pursuing the full continuum route is not an attractive path: the formulation lives on
continuum objects (measures on compact metric spaces, densities, and associated regularity/continuity assumptions), while
our satellite questions are naturally finite and simulator-driven. We therefore retain the equilibrium definition as a
useful theoretical lens, but we do not attempt to implement or rely on the full continuum variational machinery.



% References for citations used here are provided by the parent document (`docs/spec.tex`).
`.
% It intentionally contains no `\documentclass`, no preamble, and no `\begin{document}`.

\section{Scope}
We treat \textbf{planning} as the strategic placement and configuration of available
infrastructure under constraints such as interference, and \textbf{dimensioning} as the
minimum resource sizing required to satisfy per-user QoS targets under stochastic
load and resource outages.
Our three core problems of interest are interference, per-user QoS/outage, and route existence.

\section{Common Model Assumptions}
\begin{itemize}
    \item \textbf{Users} form a spatial point process on the Earth surface, within
    a latitude band.
    \item \textbf{Satellites} form a spatial point process on an orbital shell,
    gated by visibility.
    \item \textbf{Radio model} uses a received-power formulation
    $P(x,y)=P_t G_{xy}\|x-y\|^{-\gamma}$ with Rayleigh fading
    $G_{xy}\sim \mathrm{Exp}(1)$ and SINR-based capacity $C=\mathcal{W}\log_2(1+\mathrm{SINR})$.
    This follows standard stochastic-geometry modeling conventions \cite{gomez2018stochastic}.
\end{itemize}

Definitions: $P_t$ is per-satellite transmit power (W); $G_{xy}$ is a unitless
small-scale fading power gain for the link $x\leftrightarrow y$; $\|x-y\|$ is
3D link distance (m); $\gamma$ is the path-loss exponent; $\mathcal{W}$ is bandwidth (Hz);
$\mathrm{SINR}=S/(I+N)$ with $S$ the received signal power, $I$ aggregate
interference, and $N=N_0\mathcal{W}$ the noise power with $N_0$ noise spectral density (W/Hz).

\section{Planning: Interference as the Primary Objective}
We view planning as choosing configuration variables (e.g., satellite density,
altitude, power, bandwidth, and visibility constraints) so that variables of interest are
controlled (primarily intereference) while maintaining acceptable service quality.
\begin{itemize}
    \item \textbf{Interference metric}: for user at $x$ with serving satellite
    $y_0$, define $I(x)=\sum_{y\in \Phi_b\setminus\{y_0\}} P(x,y)$, that is, the aggregate power
    received from all visible satellites except the serving one.
    \item \textbf{Planning target}: enforce a region-level constraint such as
    $\mathbb{P}(\mathrm{SIR}(x)\ge\beta)\ge p_{\mathrm{sir}}$ or
    $\mathbb{E}[I(x)]\le I_{\max}$ for users in the band.
    \item \textbf{Equilibrium framing}: treat operators or satellites as agents
    whose actions affect aggregate interference, then study equilibria under
    congestion penalties.
\end{itemize}

\section{Dimensioning: Per-User QoS Under Network Outage}
Dimensioning seeks the minimum satellite intensity $N$ such that per-user QoS
targets are met with high confidence, even when the network is resource-limited.
\begin{itemize}
    \item \textbf{Per-user outage}: $O_i=1$ if a user is unserved or if
    throughput $C_i < R_{\min}$; otherwise $O_i=0$.
    \item \textbf{Outage rate}: $o = \frac{1}{|\mathcal{U}|}\sum_{i\in\mathcal{U}} O_i$.
    \item \textbf{Dimensioning target}: find the smallest $N$ such that
    $\mathbb{P}(o \le o_{\max}) \ge p_{\mathrm{target}}$ under the stated priors.
\end{itemize}

\section{Routing: Route Existence Under Time-Varying Connectivity}
We study whether an end-to-end route exists between endpoints given a time-varying visibility/connectivity graph
(ground--satellite and inter-satellite links), reporting $\mathbb{P}(\text{route exists})$ (or a certified lower bound).

\section{Planned Experiments}
\begin{itemize}
    \item Vary satellite intensity and compute interference distributions for
    latitude bands to identify feasible planning regimes.
    \item Measure per-user outage and certify feasibility over Monte Carlo trials,
    reporting the minimum $N$ that meets $(o_{\max}, p_{\mathrm{target}})$.
    \item Measure route existence (as a graph-connectivity event) under representative visibility/connectivity snapshots.
\end{itemize}



\section{Theoretical Structure}

\subsection{Topological and Metric Spaces}

\paragraph{Motivation.}
Many sets we work with in applications (such as subsets of $\mathbb{R}^d$) come with a natural notion of
``closeness'' and ``continuity''. This structure is what lets us talk about open sets, continuous costs, and Borel sets,
which in turn are the standard measurable sets used for probability measures.

\paragraph{Intuition.}
A \emph{topology} tells us which subsets of a set should be considered ``open'', generalising open intervals in
$\mathbb{R}$ and open balls in $\mathbb{R}^d$. A \emph{metric} is a concrete way to specify closeness by giving distances.
Compactness is a finiteness condition (no ``escape to infinity'') that makes many existence arguments behave well.

\paragraph{Definition.}
A \textbf{topological space} is a pair $(X,\tau)$ where $\tau\subseteq 2^X$ is a collection of subsets of $X$ (called
\emph{open sets}) such that:
$$
\emptyset\in\tau,\quad X\in\tau,\quad
\bigcup_{i\in I} U_i \in\tau\ \text{for any family }(U_i)_{i\in I}\subseteq\tau,\quad
U\cap V\in\tau\ \text{for all }U,V\in\tau.
$$

A \textbf{metric space} is a pair $(X,d)$ where $d:X\times X\to[0,\infty)$ satisfies (i) $d(x,y)=0\Leftrightarrow x=y$,
(ii) symmetry, and (iii) the triangle inequality. Every metric $d$ induces a topology by taking open sets to be unions of
open balls $B_r(x)=\{z\in X:d(z,x)<r\}$.

A metric space $(X,d)$ is \textbf{compact} if every open cover of $X$ has a finite subcover (equivalently in metric
spaces, every sequence has a convergent subsequence). A \textbf{compact metric space} is a metric space that is compact.

\paragraph{Example.}
$X=[0,1]\subset\mathbb{R}$ with the usual distance $d(x,y)=|x-y|$ is a compact metric space.
More generally, every closed and bounded subset of $\mathbb{R}^d$ with the Euclidean distance is compact.

\subsection{Measurability}

\paragraph{Motivation.}
We wish to make quantitative statements (such as probabilities of events). In order to do so, we
must ensure that the sets we are making statements about are measurable, that is, we must ensure that 
a simple real number can be assigned to their subsets in a consistent manner that correctly reflects their 
relative sizes (masses). For example, we wish to talk about probabilities of events such as 
``signal to noise and interference is below threshold with probability''.

\paragraph{Intuition.}
How do we create a framework that allows us to measure subsets of a set (e.g. events in the case of probability)? 

\paragraph{Definition.}
A \textbf{measurable space} is a pair $(X,\mathcal{F})$ where $X$ is a set and $\mathcal{F}$ is a collection of subsets
of $X$ (called ``events'') on which we will be allowed to assign sizes.
A \textbf{measure} on $(X,\mathcal{F})$ is a function $\mu:\mathcal{F}\to[0,\infty]$ such that $\mu(\emptyset)=0$ and
for any pairwise disjoint sets $A_1,A_2,\dots \in \mathcal{F}$,
$$
\mu\!\left(\bigcup_{n=1}^\infty A_n\right)=\sum_{n=1}^\infty \mu(A_n).
$$
A \textbf{probability measure} is a measure with $\mu(X)=1$.

\paragraph{Example.}
If $X=\{1,2,\dots,n\}$ is finite, a natural choice is $\mathcal{F}=2^X$ (all subsets), and the uniform probability
measure $\mu(A)=|A|/n$. Then $\mu(A)$ is the probability of the event ``the outcome lies in $A$''.

\subsection{$\sigma$-algebras}

\paragraph{Motivation.}
In probability, we want to manipulate events using the usual set operations: we negate events, take unions of events,
and take limits of event sequences. The collection of ``allowed'' events should therefore be closed under these
operations, so that we do not leave the space of events when doing algebra.

\paragraph{Intuition.}
We do \emph{not} assume every subset of $X$ is measurable: we pick a structured family of subsets that is large enough
for the geometric/physical events we care about, but stable under complements and countable unions.

\paragraph{Definition.}
This is done by equipping these sets (e.g. $X$ the set of all agent types, $Y$ the set of all actions) with a 
$\sigma$-algebra $\mathcal{F}$, which is a collection of subsets defined in such a way that assigning a real number 
to each subset in $\mathcal{F}$ is consistent with the usual rules of set operations (union, intersection, complement) 
and countable limits. In other words, we may add, subtract, and take limits while preserving the consistency of our measure.

A $\sigma$-algebra $\mathcal{F}$ on a set $X$ is a collection of subsets of $X$ such that
$$
X \in \mathcal{F}, \qquad
A \in \mathcal{F} \Rightarrow X \setminus A \in \mathcal{F}, \qquad
A_1,A_2,\dots \in \mathcal{F} \Rightarrow \bigcup_{n=1}^\infty A_n \in \mathcal{F}.
$$
These closure properties ensure that measures are consistent under negation and countable limits.

\paragraph{Example.}
When $X\subset \mathbb{R}^d$, we want
to treat ``geometric'' subsets of $X$ as legitimate events, e.g. open balls, half-spaces,
and threshold sets such as $\{x\in X : f(x)\le a\}$.

We therefore take $\mathcal{F}=\mathcal{B}(X)$, the \textbf{Borel $\sigma$-algebra} on $X$,
defined as the \emph{smallest} $\sigma$-algebra that contains all open subsets of $X$
(equivalently, the $\sigma$-algebra generated by open balls).

Open sets matter here because geometric events (``within radius $r$ of a point'', 
``inside a coverage footprint'') are naturally described using open balls and their unions, 
and we want those events to be measurable by default.

\subsection{Measurable Functions}

\paragraph{Motivation.}
We have considered measuring sets (subsets viewed as events). We now consider measuring functions
that represent physical quantities on $X$ (SINR, received power, interference, etc.), because many events
are defined by thresholding such a function.

\paragraph{Intuition.}
Let $x$ be a random location drawn according to a probability measure $\mu$ on $X$
(for example, $\mu$ could be the distribution of user locations in a latitude band).
If $f:X\to\mathbb{R}$ is a physical quantity (such as SINR), then $f(x)$ is a random real number (a random variable).

Terminology note: in this document, when we say ``distribution'' of $x$ we mean a \emph{probability measure}
$\mu\in\mathcal{P}(X)$, so $\mu(X)=1$. (More general, non-normalised measures also appear in modeling---for example,
Poisson point process \emph{intensity} measures---but probabilities are always defined using a probability measure.)

What we want to quantify are statements like ``subset of locations where the SINR is below the threshold $a$''.
This statement corresponds to the event 
$$
\{x\in X : f(x)\le a\}.
$$
Its probability is the $\mu$-mass of that set:
$$
p = \mu(\{x\in X : f(x)\le a\}).
$$

\paragraph{Definition.}
To ensure this definition is consistent (and to allow integrals like averages), we need the set
$\{x\in X : f(x)\le a\}$ to be an element of $\mathcal{F}$ for every threshold $a$.
This is exactly what it means for $f$ to be measurable.

Formally, we equip the domain with a measurable space $(X,\mathcal{F})$ and the real line with
its Borel $\sigma$-algebra $(\mathbb{R},\mathcal{B}(\mathbb{R}))$.
A function $f:X\to\mathbb{R}$ is measurable if for every Borel set $A\in\mathcal{B}(\mathbb{R})$,
the preimage $f^{-1}(A)\in\mathcal{F}$.
In practice, it is enough to check threshold sets: $f$ is measurable if for every $a\in\mathbb{R}$,
$$
\{x\in X : f(x)\le a\}\in\mathcal{F}.
$$

\paragraph{Example.}
Finally, measurability is what makes expectations meaningful: if $f$ is integrable,
the average value of the physical quantity is $\int_X f(x)\,d\mu(x)$.
When $\mathcal{F}=\mathcal{B}(X)$ and $f$ is continuous (as for most geometric/radio quantities),
$f$ is automatically measurable.

\subsection{Push-Forward Measures}
\paragraph{Motivation.}
The notion of a push-forward measure explains how a measure on $X$ (such as the distribution of locations)
induces a corresponding measure on another space through a function. In our setting this ``other space'' is often
$\mathbb{R}$ (e.g., the real-valued SINR), but the same construction applies to maps $X\to Y$ (as in optimal transport).

\paragraph{Intuition.}
Start with a random location $x\sim \mu$ on $X$. Apply a deterministic map $f$ to obtain a new value $y=f(x)$
(for example a real-valued SINR, or a transported location in $Y$). Even though $f$ is deterministic, the output is
random because the input $x$ is random; the push-forward is simply the distribution of $y$.

\paragraph{Definition.}
Formally, let $(X,\mathcal{F})$ and $(Y,\mathcal{G})$ be measurable spaces, let $\mu$ be a measure on $(X,\mathcal{F})$,
and let $f:X\to Y$ be measurable. The \textbf{push-forward} (also called \textbf{image}) measure $f_{\#}\mu$ on
$(Y,\mathcal{G})$ is defined by
$$
(f_{\#}\mu)(B) := \mu(f^{-1}(B)), \qquad B\in\mathcal{G}.
$$
If $\mu$ is a probability measure, then $f_{\#}\mu$ is again a probability measure.
This says: the mass assigned to a set $B$ of outputs equals the mass of the set of inputs that map into $B$.

\paragraph{Example.}
In the common case $Y=\mathbb{R}$ with $\mathcal{G}=\mathcal{B}(\mathbb{R})$, letting $B=(-\infty,a]$ gives
$$
(f_{\#}\mu)((-\infty,a]) = \mu(\{x\in X : f(x)\le a\}),
$$
which is exactly the probability statement we use for threshold events.

\subsubsection{Push-Forward in the Optimal Transport Problem}
\paragraph{Motivation.}
In optimal transport, we want to \emph{move mass} from a source distribution $\mu$ on a space $X$ (a ``pile'')
to a target distribution $\nu$ on a space $Y$ (a ``hole''), while minimizing a cost (the objective).

\paragraph{Intuition.}
A transport map $T:X\to Y$ tells us where each location $x$ sends its mass. The requirement that ``the pile becomes
the hole'' is exactly that transporting $\mu$ through $T$ yields the target distribution $\nu$.

\paragraph{Definition.}
We say that $T$ \emph{pushes $\mu$ forward to $\nu$} if
$$
T_{\#}\mu=\nu,
$$
meaning that for every measurable set $B\subseteq Y$,
$$
\nu(B)=\mu(T^{-1}(B)).
$$
This is exactly the push-forward definition from the previous subsection, with the specific map $f$ renamed to $T$
(to emphasize that it is a transport map) and the codomain being $Y$ rather than $\mathbb{R}$.

\textbf{Kantorovich formulation.} The Monge formulation above restricts us to a \emph{deterministic} assignment:
each $x$ sends all its mass to a single point $T(x)$. This can be too restrictive (for example, if a point of $\mu$
must supply multiple points in $\nu$) and it leads to a nonconvex optimization problem.

The Kantorovich idea is to relax ``one destination per source point'' into a \emph{transport plan} that can split mass:
instead of choosing a map $T$, we choose a joint measure $\gamma$ on $X\times Y$ whose mass at $(x,y)$ represents
how much mass is transported from $x$ to $y$.

In optimisation language, the \textbf{decision variable} (also called the \textbf{optimisation variable}) is the
mathematical object we are free to choose in order to minimise the objective under constraints.
In the Kantorovich formulation, the decision variable is the joint measure $\gamma$ on $X\times Y$.

To represent a valid transport plan from $\mu$ to $\nu$, $\gamma$ should have marginals $\mu$ and $\nu$.
Equivalently, the constraints that $\gamma$ has marginals equal to $\mu$ and $\nu$ can be written as push-forward relations
$$
(\pi_X)_{\#}\gamma=\mu,\qquad (\pi_Y)_{\#}\gamma=\nu,
$$
where $\pi_X(x,y)=x$ and $\pi_Y(x,y)=y$ are the coordinate projections.
These are exactly push-forwards of the same form $T_{\#}(\cdot)$, with the ``map'' $T$ being instantiated as
$T=\pi_X$ and $T=\pi_Y$, respectively.

To make the origin and endpoint explicit: $\gamma$ is a measure on the \emph{product space} $X\times Y$.
The map $\pi_X:X\times Y\to X$ takes a pair $(x,y)$ and returns its $X$-coordinate $x$, so $(\pi_X)_{\#}\gamma$
is a measure \emph{on $X$}. Likewise, $\pi_Y:X\times Y\to Y$ returns the $Y$-coordinate, so $(\pi_Y)_{\#}\gamma$
is a measure \emph{on $Y$}.

\paragraph{Example.}
Let $X=[0,1]$, $Y=[0,2]$, and let $\mu$ be the uniform distribution on $[0,1]$. Define $T(x)=2x$.
Then $T_{\#}\mu$ is the uniform distribution on $[0,2]$ (the interval is stretched, so density halves), i.e.\ $\nu=T_{\#}\mu$.

Kantorovich-only example (mass splitting). Let $X=\{0\}$ and $Y=\{-1,+1\}$. Let $\mu=\delta_0$ (all mass at $0$) and
let $\nu=\tfrac12\delta_{-1}+\tfrac12\delta_{+1}$. There is no map $T:X\to Y$ such that $T_{\#}\mu=\nu$, because any map
must send all mass to a single point in $Y$. But there is a transport plan $\gamma$ on $X\times Y$ defined by
$$
\gamma(0,-1)=\tfrac12,\qquad \gamma(0,+1)=\tfrac12,
$$
which satisfies $(\pi_X)_{\#}\gamma=\mu$ and $(\pi_Y)_{\#}\gamma=\nu$.

\section{Optimal Transport and Cournot-Nash Equilibria}
This section summarises the background and key ideas in \cite{blanchet2016optimal} in a format aligned with our study.

We are interested in settings with many agents, each attempting to optimise for itself. What sorts of equilibria arise?
What do these equilibria look like? When do they exist, and when are they unique?

For example, consider a large number of users (agents) choosing which satellite (action) to connect to. 
Each user's utility depends on which satellite it chooses, but also on how many other users choose the 
same satellite (congestion externality). We wish to characterise the equilibria of this system, so we may 
understand not only how users behave but also how to design the system (e.g., satellite parameters) to 
achieve desirable outcomes.

Each user is solving an individual optimisation problem, but the problems are coupled through the congestion effects.
We need to characterise jointly the optimisation problem as well as the equilibrium conditions. It would not suffice 
to simply solve each user's optimisation problem in isolation, or to show that an equilibrium exists. It is essential 
to simultaneously characterise individual optimisation strategies and the resulting equilibrium.

In game-theoretic terms, we are interested in \textbf{non-atomic games with externalities}, where each agent is infinitesimally
small (non-atomic) and the cost incurred by each agent depends on the aggregate behaviour of all agents (externalities).

\subsection{Problem Class, Background, and Gap}

\cite{blanchet2016optimal} studies non-atomic games with externalities and gives classes of models where Cournot--Nash
equilibria are not only shown to exist, but can also be characterised in a way that supports uniqueness results and
computation. Their stated gap is that, while abstract existence results for Cournot--Nash equilibria are well known, they
do not typically yield a tractable characterisation (or numerical scheme) for the equilibrium distribution.

\paragraph{Where do externalities and congestion come from?}
In many applications, an agent's cost depends on what \emph{others} do through an aggregate quantity. For instance, in a
wireless association model, if many users choose the same satellite (or the same beam), then service quality degrades;
in road choice, if many drivers choose the same road, travel time increases. These effects are called \emph{externalities}
because an agent's cost depends on the aggregate behaviour of the population, not only on the agent's own action.
The paper models such effects through an action distribution $\nu$ and an externality term $V[\nu](y)$ that modifies the
cost of playing action $y$.

However, once we introduce such externalities (especially congestion-type effects), there is a key technical
obstruction: the dependence of costs on the population distribution can fail to have the regularity properties assumed
by standard equilibrium existence proofs.

\paragraph{What is the ``obstruction''?}
Here, an \emph{obstruction} simply means a technical barrier that prevents us from applying the most common existence
proof techniques ``out of the box''. Concretely, many existence proofs for equilibria rely on a fixed-point argument and
therefore need a basic stability property: if the population action distribution $\nu$ changes a little, then the cost
landscape faced by agents should also change a little. Without such stability, it becomes hard to even define a
well-behaved ``best-response of the population'' map, let alone prove it has a fixed point.

Congestion effects are often \emph{local}: the cost at action $y$ depends on how crowded that action is. A common model is
$V[\nu](y)=f(\nu(y))$, where $\nu(y)$ denotes a density. The difficulty is that densities are fragile: two action
distributions can look almost the same when you only look at coarse averages (e.g.\ averages over regions), yet have very
different pointwise densities because mass can concentrate into narrow spikes or oscillate rapidly. If $V[\nu](y)$ depends
directly on $\nu(y)$, these small ``distribution-level'' changes can produce large changes in the local congestion cost,
so the regularity assumptions behind standard fixed-point existence arguments fail.

To make these statements precise, we first fix the paper's equilibrium notation (Subsection~\ref{sec:cn_equilibrium}) and then explain
the relevant notion of ``two distributions being close'' (weak-$*$ topology) and why local congestion breaks it
(Subsection~\ref{sec:cn_weakstar_gap}).

The key contribution in \cite{blanchet2016optimal} is to exploit an optimal-transport structure (and a variational formulation in potential
cases) to turn equilibrium computation into an optimisation problem over distributions.



\subsection{Equilibrium}\label{sec:cn_equilibrium}
In a finite-player game, one can describe an equilibrium by listing each player's action (or mixed strategy).
In a \emph{non-atomic} game with a continuum of agents, this is not the right object: we do not label agents individually,
and a single agent has negligible influence on the aggregate outcome. Instead, we describe equilibria at the population
level using probability measures.

\subsubsection{Spaces and Probability Measures}
In the paper's notation, $X$ is the \emph{type space} (agent characteristics) and $Y$ is the \emph{action space}.
We write $\mathcal{P}(X)$ for the set of \emph{Borel} probability measures on $X$ and similarly $\mathcal{P}(Y)$.
Here ``Borel'' means the measures are defined on the Borel $\sigma$-algebra generated by the topology on $X$ (and $Y$),
which is the standard choice when $X$ and $Y$ are compact metric spaces. (The ``continuum of agents'' aspect is captured
by describing populations via measures.)
The population of types is given exogenously by a probability measure $\mu\in\mathcal{P}(X)$.

\subsubsection{Equilibrium as a Joint Distribution}
What matters is the population-level assignment of actions to types: for each type $x$, what fraction of agents choose
each action $y$. This assignment is represented by a joint probability measure on $X\times Y$.

An equilibrium is described by a joint probability measure $\gamma\in\mathcal{P}(X\times Y)$, where (informally)
$\gamma(A\times B)$ is the probability that an agent has type in $A\subseteq X$ and chooses an action in $B\subseteq Y$.
The first marginal of $\gamma$ is the induced distribution of types and must equal $\mu$.
The second marginal of $\gamma$ is the induced distribution of actions, denoted $\nu\in\mathcal{P}(Y)$.
Using the coordinate projections $\pi_X(x,y)=x$ and $\pi_Y(x,y)=y$, these marginal constraints can be written as
push-forward relations:
$$
(\pi_X)_{\#}\gamma=\mu,\qquad \nu=(\pi_Y)_{\#}\gamma.
$$
If the equilibrium is \emph{pure}, then there is a map $T:X\to Y$ such that $\gamma$ is carried by the graph of $T$,
i.e.\ $\gamma=(\mathrm{id},T)_{\#}\mu$.

\subsubsection{Costs and Externalities}
Costs are additively separable: an $x$-type agent who plays action $y$ incurs a cost of the form
$$
\Phi(x,y,\nu)=c(x,y)+V[\nu](y),
$$
where $c(x,y)$ is an idiosyncratic term (distance, mismatch, etc.), and the map $\nu\mapsto V[\nu]$ captures
externalities induced by the aggregate action distribution $\nu$.

\subsubsection{Reference Measure, Densities, and Local Congestion}
To speak about ``local density'' at an action $y$, we need $\nu$ to admit a density with respect to some reference
measure on $Y$.
Fix a reference measure $m_0$ on $Y$ and restrict to action distributions $\nu$ that are absolutely continuous with
respect to $m_0$ (written $\nu\ll m_0$). In that case, $\nu$ admits a density $\frac{d\nu}{dm_0}$, and with a slight abuse
of notation we write $\nu(y)$ for this density value at $y$.
Local congestion costs take the form
$$
V[\nu](y)=f(\nu(y)),
$$
so the externality cost at $y$ depends pointwise on how concentrated the action distribution is at $y$.

\subsubsection{Cournot--Nash Equilibrium Condition (Formal Definition)}\label{sec:cn_equilibrium_definition}
\paragraph{Motivation.}
So far, $\gamma$ is just a joint distribution that matches the exogenous type distribution $\mu$ and induces an action
distribution $\nu$. To call $\gamma$ an \emph{equilibrium}, we must encode the behavioural requirement that agents choose
cost-minimising actions \emph{given} the aggregate action distribution $\nu$ that results from everyone's choices.

\paragraph{Intuition.}
Fix a candidate aggregate action distribution $\nu$. Each type $x$ then faces a deterministic optimisation problem:
choose $y\in Y$ to minimise $c(x,y)+V[\nu](y)$. An equilibrium is a self-consistent joint distribution $\gamma$ such that
(i) its second marginal is exactly the $\nu$ used in the optimisation, and (ii) $\gamma$ puts mass only on pairs $(x,y)$
where $y$ is a best response for type $x$.

\paragraph{Definition.}
Let $\gamma\in\mathcal{P}(X\times Y)$ and let $\nu=(\pi_Y)_{\#}\gamma$ be its second marginal.
We say that $\gamma$ is a \textbf{Cournot--Nash equilibrium} (in distributional form) if
\begin{enumerate}[label=(\roman*)]
    \item $(\pi_X)_{\#}\gamma=\mu$ (types are distributed according to $\mu$), and
    \item $\gamma$ is concentrated on best responses given $\nu$, meaning
    $$
    \gamma\bigl(\{(x,y)\in X\times Y:\ c(x,y)+V[\nu](y)=\min_{z\in Y}\bigl(c(x,z)+V[\nu](z)\bigr)\}\bigr)=1.
    $$

    In probabilistic terms: let $(\Omega,\mathcal{F},\mathbb{P})$ be a probability space and let $(X_{\Omega},Y_{\Omega})$
    be an $(X\times Y)$-valued random variable with law $\gamma$ (so $\mathbb{P}[(X_{\Omega},Y_{\Omega})\in A]=\gamma(A)$ for
    every measurable $A\subseteq X\times Y$). Define the event
    $$
    E := \Bigl\{c(X_{\Omega},Y_{\Omega})+V[\nu](Y_{\Omega})=\min_{z\in Y}\bigl(c(X_{\Omega},z)+V[\nu](z)\bigr)\Bigr\}.
    $$
    Then the condition $\gamma(\cdot)=1$ above is exactly the statement $\mathbb{P}(E)=1$: with probability one, the sampled
    action $Y_{\Omega}$ is a best response for the sampled type $X_{\Omega}$ given the aggregate distribution $\nu$.

\end{enumerate}


\paragraph{Intuition.}
This expresses a ``best response for almost every agent'' condition in a non-atomic (continuum) model: we do \emph{not}
quantify over individually labelled agents. Instead, we quantify over the population measure $\gamma$: the set of
non-best-response pairs has $\gamma$-mass $0$, i.e.\ if we pick an agent at random from the population (sample
$(x,y)\sim\gamma$), then with probability $1$ the realised action $y$ is a best response for the realised type $x$.

Equivalently, if we define the best-response set $B_{\nu}\subseteq X\times Y$ by the bracketed condition inside
$\gamma(\cdot)$ above, then $\gamma(B_{\nu})=1$ and the event $E$ is exactly $\{(X_{\Omega},Y_{\Omega})\in B_{\nu}\}$.

\paragraph{Alternative Formulation.}

Equivalently, define the \emph{value function}
$$
\varphi(x):=\min_{z\in Y}\bigl(c(x,z)+V[\nu](z)\bigr),
$$
then the equilibrium condition above can be written as the pair of inequalities
$$
c(x,y)+V[\nu](y)\ge \varphi(x)\quad\text{for all }x\in X\text{ and (where defined) }y\in Y,
$$
with equality holding $\gamma$-almost everywhere (i.e.\ for $\gamma$-almost every pair $(x,y)$).

\paragraph{Example.}
In a satellite-association toy model, let $x$ be a user location and let $y$ index a satellite/beam choice. The term
$c(x,y)$ can represent a geometry-based penalty (e.g.\ path loss), and $V[\nu](y)$ can represent a congestion penalty for
how many users choose satellite/beam $y$ (captured by the aggregate action distribution $\nu$). The equilibrium
condition says that, under the self-consistent distribution of choices, almost every user is assigned only to a choice
that minimises its own cost given the congestion created by all users.


\paragraph{Motivation.}
The equilibrium definition in Subsubsection~\ref{sec:cn_equilibrium_definition} is a self-consistent best-response
condition on a joint distribution $\gamma$.
This definition is conceptually clear, but it does not yet tell us how to \emph{compute} an equilibrium, or how to obtain
uniqueness and stability properties.

\paragraph{Intuition.}
The paper's key step is to identify an important class of models where the equilibrium conditions are exactly the
first-order optimality conditions of a single global objective. In that case, equilibrium computation becomes an
optimisation problem over distributions. Optimal transport enters through the transport cost term, while congestion and
interactions enter through an ``energy'' functional of the action distribution.

\paragraph{Definition.}
In the \emph{potential} (variational) case, the externality map $\nu\mapsto V[\nu]$ is assumed to be the first variation
of some functional $\mathcal{E}$ defined on action distributions:
$$
V[\nu] \;=\; \frac{\delta \mathcal{E}}{\delta \nu}.
$$
Under this assumption, equilibrium action distributions $\nu$ can be characterised as minimisers of an OT-structured
functional of the form
$$
\nu \in \arg\min_{\tilde\nu}\ \mathcal{W}_c(\mu,\tilde\nu) + \mathcal{E}[\tilde\nu],
$$
and an equilibrium coupling $\gamma$ is then obtained as an optimal transport plan between $\mu$ and that minimising
$\nu$.

\paragraph{Example.}
A standard ``good'' potential form used in the paper is
$$
\mathcal{E}[\nu] \;=\; \int_Y F(y,\nu(y))\,dm_0(y) \;+\; \frac12\iint_{Y\times Y} \phi(y,z)\,d\nu(y)\,d\nu(z),
$$
where the first term is a local congestion energy (with $F'(\cdot)=f(\cdot)$) and the second term is a symmetric pairwise
interaction energy (symmetry is what makes it a potential).

\subsection{Why Variational Structure Matters}\label{sec:cn_variational_importance}

\paragraph{Motivation.}
The ability to write equilibrium computation as a global minimisation problem is the reason this approach becomes
practical. Without a variational structure, we typically lose both tractability and the clean uniqueness/stability
statements that come from convexity (or convexity along transport geodesics).

\paragraph{Intuition.}
For general Nash equilibria, there is no reason that ``the equilibrium'' should be the minimiser of a single scalar
objective: different players optimise different objectives, and their incentives can create cycles and rotational
behaviour. A variational (potential) structure is exactly the condition that makes the equilibrium conditions integrable
into a single energy landscape: equilibria are stationary points of that landscape, and in the convex cases they are
global minimisers.

In the Cournot--Nash setting considered by \cite{blanchet2016optimal}, the variational structure is what turns a
definition-by-best-response (Subsubsection~\ref{sec:cn_equilibrium_definition}) into a computable characterisation
(Subsection~\ref{sec:cn_variational_formulation}): solve a minimisation problem for the action distribution $\nu$, then
recover the coupling $\gamma$ using optimal transport.

\paragraph{Definition.}
We will refer to the model as \textbf{variational} (or \textbf{potential}) if there exists a functional
$\mathcal{E}$ on action distributions such that $V[\nu]=\delta\mathcal{E}/\delta\nu$.
In this case, the equilibrium conditions coincide with the first-order optimality conditions of the global objective
$\mathcal{W}_c(\mu,\nu)+\mathcal{E}[\nu]$.
If, in addition, this objective is (geodesically) convex in a suitable sense, then the equilibrium is unique and stable
to perturbations of the inputs $(\mu,c,\mathcal{E})$.

\paragraph{Example.}
Not every game has this structure. In finite games such as Rock--Paper--Scissors or Matching Pennies, incentives create
cyclic dominance, and there is no single potential function whose value strictly improves along unilateral best-response
moves. In continuous-action games, strong asymmetric cross-effects can lead to non-monotone best-response operators (the
system behaves more like a rotation than a gradient flow). The relevance here is conceptual: the Blanchet--Carlier
approach relies on being able to convert equilibrium computation into minimisation, which is precisely what fails when a
game is not variational.

\subsection{Practical Takeaway: Variational Characterisation vs.\ Continuum Complexity}
\label{sec:cn_weakstar_gap}
\label{sec:cn_variational_formulation}
The main technical payoff in \cite{blanchet2016optimal} is that, in the \emph{potential} (variational) case, equilibrium
computation can be reframed as a single global optimisation problem (an optimal-transport assignment term plus an energy
term in the action distribution). This provides a principled \emph{characterisation} of equilibrium and, in favourable
convex cases, supports uniqueness and stable numerical schemes.

For this study, however, pursuing the full continuum route is not an attractive path: the formulation lives on
continuum objects (measures on compact metric spaces, densities, and associated regularity/continuity assumptions), while
our satellite questions are naturally finite and simulator-driven. We therefore retain the equilibrium definition as a
useful theoretical lens, but we do not attempt to implement or rely on the full continuum variational machinery.



% References for citations used here are provided by the parent document (`docs/spec.tex`).
